\providecommand{\keyentry}[7]{%
  \par\addvspace{1.6ex}\noindent
  \begin{minipage}[t]{4.6em}\raggedright\setlength{\parskip}{0pt}%
    \ifnum#1=3 \small\fi
    \makebox[1.1em][l]{\ifcase#1\or ●\or ○\or\or ～\fi}#2\par
    {\scriptsize\textcolor{dimink}{#3}}\end{minipage}%
  \begin{minipage}[t]{\dimexpr\linewidth-4.6em\relax}\raggedright\setlength{\parskip}{0.4ex}%
    \ifnum#1=3 \small\fi
    \if\relax\detokenize{#4}\relax\else\noindent{\small\textcolor{dimink}{#4}}\par\fi
    {\setstretch{1.2}\noindent #6\par}\nopagebreak
    {\ifnum#1=3 \footnotesize\else\small\fi\setstretch{1.3}\noindent\textcolor{dimink}{#5}\par}\nopagebreak
    {\ifnum#1=3 \footnotesize\else\small\fi\setstretch{1.3}\noindent\textcolor{dimink}{理由}\quad #7\par}\end{minipage}\par}

% 本檔由 gen_keymarks_qa.py 自翻譯倉 Notizen/導讀標記.tsv 產生，勿手改。
\qaentry{第二部之四}{重點句的分級與論證骨架}

\begin{qaquestion}
導讀時哪些句子要逐字講出？依什麼標準分級？哪幾句連起來就是全文的論證？
\end{qaquestion}

本文共標 253 句，分四級。前三級依各句在論證中的作用分；第四級收作者生動的表述（比喻、典故、反詰），不論它在論證中的位置。編號以 K 起首，本書各則引到某句時即用此編號；中譯取自譯本。各句依原文順序列於第二部之五。

\begin{enumerate}
\item 一級 10 句：全文的主張、攻擊標的、論證的條件與正面主張、兩類理由各自的結論、論證的範圍及其限度，以及結論兩句。依第二部之四所列順序講完，就是全文的論證；任刪一句，論證就缺一步。
\item 二級 79 句：各部分的主旨與作者回應常見反駁之處，講到該部分時逐字講出。
\item 三級 94 句：定義、舉例、預告與過渡，時間不足時一句帶過。
\item 四級 70 句：作者生動的表述，講到該處時可以照念。
\end{enumerate}

各部分的分布如下表。

\begin{longtable}{@{}L{0.26\linewidth}L{0.10\linewidth}L{0.16\linewidth}L{0.16\linewidth}L{0.16\linewidth}@{}}
\toprule
部分 & 一級 & 二級 & 三級 & 四級 \\
\midrule
\endhead
摘要 & \small K01 & — & \small K02 & — \\
導論 & — & \small K05 K50 K182 & \small K03 K04 K48 K49 K82 K83 K84 K85 K86 K87 K88 K184 & \small K183 K185 K199 K200 K201 K202 \\
壹、違憲審查的界定 & \small K06 & \small K07 K08 & \small K51 K52 K53 K54 K55 K56 K89 K149 K186 & \small K203 K204 \\
貳、四項假設 & \small K09 K10 & \small K11 K14 K16 K17 K19 K58 K59 K90 K91 K92 K94 K95 K96 K97 K151 & \small K12 K13 K15 K18 K57 K93 K98 K150 K152 K153 K154 & \small K205 K206 K207 K208 K209 K210 \\
參、論證的形式 & \small K22 K23 & \small K20 K21 K62 K63 K65 K100 K101 K102 K103 K156 & \small K60 K61 K64 K99 K104 K105 K155 K157 K188 & \small K187 K189 K211 K212 \\
肆、結果面的理由 & — & \small K26 K27 K28 K29 K69 K70 K108 K110 K113 K158 K192 & \small K24 K25 K31 K66 K67 K68 K106 K107 K109 K111 K112 K114 K115 K159 K160 K163 & \small K30 K161 K162 K190 K191 K213 K214 K215 K216 K217 K218 K219 K220 K221 K222 K223 K224 K225 K226 K227 K228 K229 K230 \\
伍、程序面的理由 & — & \small K33 K34 K35 K36 K37 K73 K74 K75 K116 K118 K120 K125 K126 K127 K128 K166 K168 & \small K32 K71 K72 K117 K119 K121 K122 K123 K124 K147 K148 K164 K165 K167 & \small K193 K194 K231 K232 K233 K234 K235 K236 K237 K238 K239 K240 K241 K242 \\
陸、多數暴政 & — & \small K39 K40 K41 K76 K77 K130 K131 K134 K135 K139 K170 K171 & \small K38 K129 K132 K133 K136 K137 K138 K169 K172 K173 K174 & \small K195 K196 K243 K244 K245 K246 K247 K248 \\
柒、非核心情形 & \small K42 K43 & \small K44 K79 K81 K140 K144 K176 K178 K181 & \small K45 K78 K80 K141 K142 K143 K145 K175 K177 K179 K180 & \small K197 K198 K249 K250 K251 K252 K253 \\
結論 & \small K46 K47 & \small K146 & — & — \\
合計 & 10 & 79 & 94 & 70 \\
\bottomrule
\end{longtable}

一級 10 句在原文中的先後，即論證的步驟：

\begin{longtable}{@{}L{0.22\linewidth}L{0.12\linewidth}L{0.20\linewidth}L{0.36\linewidth}@{}}
\toprule
步驟 & 編號 & 刊本頁 & 部分 \\
\midrule
\endhead
主張 & K01 & 1346 & 摘要 \\
標的 & K06 & 1354 & 壹、違憲審查的界定 \\
條件 & K09 & 1360 & 貳、四項假設 \\
正面主張 & K10 & 1360 & 貳、四項假設 \\
結果面 & K22 & 1375 & 參、論證的形式 \\
程序面 & K23 & 1375 & 參、論證的形式 \\
範圍 & K42 & 1402 & 柒、非核心情形 \\
（同上） & K43 & 1402 & 柒、非核心情形 \\
結論一 & K46 & 1406 & 結論 \\
結論二 & K47 & 1406 & 結論 \\
\bottomrule
\end{longtable}


\qaentry{第二部之五}{各部分的重點句}

\begin{qaquestion}
各句在各部分的哪個位置？各為何入選？
\end{qaquestion}

依原文的部分與段落先後排列。左欄為層級符號、編號與刊本頁：● 一級，○ 二級，～ 四級，無符號者為三級，三級並以小一號的字排。一級句另於右欄首行標出它在論證中的步驟。

\subsection*{摘要}
\keyentry{1}{K01}{1346}{一級・主張}{本文從兩個主要理由批評違憲審查。其一，沒有理由認為權利在這個制度下會比在民主立法機關下受到更好的保護。其二，撇開它產生的結果，違憲審查在民主上欠\mbox{缺正當性。}}{The Essay criticizes judicial review on two main grounds. First, it argues that there is no reason to suppose that rights are better protected by this practice than they would be by democratic legislatures. Second, it argues that, quite apart from the outcomes it generates, judicial review is democratically illegitimate.}{全文的兩個理由在此並列：結果面無從認定法院較能保障權利，程序面欠缺民主正當性。第肆、伍兩部分依序展開這\mbox{兩個理由。}}
\keyentry{3}{K02}{1346}{}{然而，反對違憲審查的論據並非絕對，也\mbox{非無條件。}}{However, the case against judicial review is not absolute or unconditional.}{論證附條件的預告，正式的表述在\mbox{ K09。}}
\subsection*{導論}
\keyentry{3}{K82}{1348}{}{不過他們認為，這些代價常被誇大，或遭到曲解；在他們看來，民主程序絕非完美，而且無論如何，問題涉及多數暴政時，以民主為理由的反對本身\mbox{就站不住。}}{But, they say, these costs are often exaggerated or mischaracterized. The democratic process is hardly perfect and, in any case, the democratic objection is itself problematic when what is at stake is the tyranny of the majority.}{辯護者的立場：違憲審查的代價常被誇大，多數暴政之下以民主為由的反對站不住。本段以 Goodridge、Lawrence、Roe、Brown 等判決為其成果。第陸部分回應多數\mbox{暴政之說。}}
\keyentry{3}{K03}{1348}{}{這差不多是我對違憲審查要說的最後\mbox{一句好話。}}{That is almost the last good thing I shall say about judicial review.}{導論前段肯定 Goodridge 一類判決的成果，此句起\mbox{轉入批評。}}
\keyentry{3}{K83}{1349}{}{Alexander Bickel 用一句知名的話概括了這個問題：「對抗多\mbox{數的難題」。}}{Alexander Bickel summed up the issue in the well-known phrase, ``the counter-majoritarian difficulty.''}{既有反對論述的代表用語。Bickel 的緩和辦法與其限度見下一\mbox{段的引文。}}
\keyentry{4}{K199}{1349}{四級・生動表述}{這個體系還有更多的複雜與費解之處，現代政治學以令人欽佩的巧思與勤勉探索了它們，其中有些經它一再增衍，其多產已遠非發現者的熱忱所能解釋——但這些複雜之處，無一能改變以下根本事實：在美國的民主裡，違憲審查是\mbox{脫軌制度。}}{nothing in the further complexities and perplexities of the system, which modern political science has explored with admirable and ingenious industry, and some of which it has tended to multiply with a fertility that passes the mere zeal of the discoverer—nothing in these complexities can alter the essential reality that judicial review is a deviant institution in the American democracy.}{Bickel 以反諷筆法揶揄政治學「多產」得超出發現者的熱忱，收在「脫軌制度」一語；導讀時可接 K83 的「反多數困境」，說明這個詞的出\mbox{處與語氣。}}
\keyentry{3}{K48}{1349}{}{在不允許以這種方式使法律失效的國家，人民自己可以透過通常的立法程序終局地決定，要不要允許墮胎、優惠性差別待遇（affirmative action）\mbox{、學券制}}{In countries that do not allow legislation to be invalidated in this way, the people themselves can decide finally, by ordinary legislative procedures, whether they want to permit abortion, affirmative action, school vouchers, or gay marriage.}{沒有違憲審查的國家，權利爭點由人民經通常立法程序終局決定。同段舉 1960 年代英國國會的墮胎、同性性行為與死刑辯論為例，下一段轉到美國由法官定\mbox{奪的情形。}}
\keyentry{2}{K182}{1349--1350}{}{那些辯論的品質（以及加拿大、澳洲、紐西蘭與其他地方類似的辯論）使「議員沒有能力負責任地處理這類議題」的說法\mbox{顯得荒謬}}{The quality of those debates (and similar debates in Canada, Australia, New Zealand, and elsewhere) make nonsense of the claim that legislators are incapable of addressing such issues responsibly}{英國國會 1960 年代就墮胎、同性性行為、死刑的辯論，作者據以反駁「議員無法負責任地處理權利議題」；第肆部分 1966 年墮胎法案辯論與 Roe 案的對比\mbox{由此預告。}}
\keyentry{3}{K84}{1350}{}{如果有人不同意立法的解決方案，決定把事情提到法院，最後勝出的是法\mbox{官的見解。}}{If someone who disagrees with the legislative resolution decides to bring the matter before a court, the view that finally prevails will be that of the judges.}{美國的情形，與 K48 所述由人民經立法程序終局決定的國家對照；同段引 Dworkin 承認大法官的洞見並不\mbox{特別出眾。}}
\keyentry{3}{K85}{1350}{}{近年出現了一些攻擊美國違憲\mbox{審查的書。}}{In recent years, a number of books have appeared attacking judicial review in America.}{美國爭論的現況：自由派原本支持違憲審查，Rehnquist 法院時期轉而出現自由派的反對；Marbury 案兩百週年與 Brown 案五十週年也引出熱\mbox{烈的辯護。}}
\keyentry{4}{K200}{1350}{四級・生動表述}{於是戰線已經劃定，攻防的手法為人熟悉，雙方的立場也早已\mbox{眾所周知。}}{So the battle lines are drawn, the maneuvering is familiar, and the positions on both sides are well understood.}{把違憲審查之爭寫成戰線已定的陣地戰；導讀時用來交代作者為何還要\mbox{再寫一篇。}}
\keyentry{4}{K201}{1351}{四級・生動表述}{我想做的，是找出一個反對違憲審查的核心論證，這個論證既獨立於違憲審查在歷史上的各種表現，也獨立於關於它特定效果的問題，亦即它所作的判決（有好有壞），以及這些判決帶給人的心\mbox{碎與肯認。}}{What I want to do is identify a core argument against judicial review that is independent of both its historical manifestations and questions about its particular effects—the decisions (good and bad) that it has yielded, the heartbreaks and affirmations it has handed down.}{「心碎與肯認」把判決的效果寫成當事人的情緒；導讀時點出本文刻意撇開的正是判決的\mbox{這些效果。}}
\keyentry{4}{K202}{1351}{四級・生動表述}{如 Frank Michelman 在 The People Themselves 封底的推薦語所說，Kramer 的歷史「為違憲審查與人民憲政主義之辯論的骨架添\mbox{上了血肉」。}}{As Frank Michelman says in his blurb on the back cover of The People Themselves, Kramer's history ``puts flesh on the bones of debates over judicial review and popular constitutionalism.''}{骨架與血肉的比喻由此起，下一句 K04 作者反過來要把血肉剝掉、熬到只剩骨架；兩\mbox{句宜連讀。}}
\keyentry{3}{K04}{1351}{}{但我想剝掉一些血肉，把規範性的論證熬到只剩骨架，好讓我們直視違憲審查，看清它以什\mbox{麼為前提。}}{But I want to take off some of the flesh and boil down the normative argument to its bare bones so that we can look directly at judicial review and see what it is premised on.}{本文的方法：撇開個案與各國歷史，只留規範論證的骨架。摘要第一段與第參部分都依此設定\mbox{討論範圍。}}
\keyentry{4}{K183}{1351}{四級・生動表述}{Charles Black 曾說，對違憲審查的反對，實際上往往「時有時無」：人們因為自己珍視的那幾件判決（例如 Brown 或 Roe）而支持違憲審查，只有當它導致自己痛惡的結果\mbox{時才反對。}}{Charles Black once remarked that, in practice, opposition to judicial review tends to be ``a sometime thing,'' with people supporting it for the few cases they cherish (like Brown or Roe) and opposing it only when it leads to outcomes they deplore.}{Charles Black「時有時無」之說：對違憲審查的態度隨判決\mbox{結果而定。}}
\keyentry{3}{K184}{1351}{}{這在墮胎權的辯論裡最為明顯：主張選擇權（pro-choice）的論者驚慌之下，連考慮反對違憲審查的論據都不肯，唯恐此舉助長另一些人的聲勢：這些人認為 Roe v. Wade 無理侵犯了保守派議\mbox{員的權利。}}{This is most notably true in the debate over abortion rights, in which there is a panic-stricken refusal among pro-choice advocates to even consider the case against judicial review for fear this will give comfort and encouragement to those who regard Roe v. Wade as an unwarranted intrusion on the rights of conservative legislators.}{支持選擇權的一方擔心動搖 Roe，連反對違憲審查的論證都不願考慮；作者說明為何把論證抽離\mbox{個案結果。}}
\keyentry{3}{K49}{1352}{}{但即使如此，仍值得究明這種辯護是否觸及問題的核心，還是只應視為一項例外的理由：反對該制度的規範論證在多數情況下會具有說服力，這項理由使我們在這些情形下暫不依循該論\mbox{證的走向。}}{But even if that is so, it is worth figuring out whether that sort of defense goes to the heart of the matter, or whether it should be regarded instead as an exceptional reason to refrain from following the tendency of what, in most circumstances, would be a compelling normative argument against the practice.}{以防範立法弊病為由支持違憲審查的辯護，作者定位為例外的理由；第柒部分的非核心情形處理\mbox{這類例外。}}
\keyentry{4}{K185}{1352}{四級・生動表述}{例如，英國人辯論其法官審查立法的權限（相對有限）時，並不特別在意 1805 年共和黨人對聯邦黨人說了什麼，也不在意 Brown v. Board of Education 留\mbox{下了什麼。}}{For example, when the British debate the relatively limited powers their judges have to review legislation, they are not particularly interested in what the Republicans said to the Federalists in 1805 or in the legacy of Brown v. Board of Education.}{以英國為例，說明論證要抽離美國的歷史與判例；接\mbox{ K86。}}
\keyentry{3}{K86}{1352}{}{需要的是某種一般性的理解，不受各個社會的文化、歷史與政治關\mbox{切所沾染。}}{What is needed is some general understanding, uncontaminated by the cultural, historical, and political preoccupations of each society.}{撇開美國歷史的第二個理由：英國等國也在辯論違憲審查，K04 的骨架式論證因此要能脫離美\mbox{國的脈絡。}}
\keyentry{3}{K87}{1352}{}{但我在討論違憲審查時，把它與法理學和政治哲學的其他問題糾纏\mbox{在了一起。}}{But I have managed to discuss judicial review in a way that embroils it with other issues in jurisprudence and political philosophy.}{作者交代本文與其舊作的關係：舊作把違憲審查與其他問題混在一起談，本文要把基本的反對理由單\mbox{獨講清楚。}}
\keyentry{2}{K50}{1353}{}{本文將主張，違憲審查在兩條戰線上都易受攻擊。它並不像人們常說的那樣，提供一種方式讓社會在公民對權利有分歧時清楚聚焦於真正的爭點；相反，它用先例、文本與解釋這些枝節問題分散了他們\mbox{的注意力。}}{In this Essay, I shall argue that judicial review is vulnerable to attack on two fronts. It does not, as is often claimed, provide a way for a society to focus clearly on the real issues at stake when citizens disagree about rights; on the contrary, it distracts them with side-issues about precedent, texts, and interpretation.}{兩條戰線的第一條：違憲審查以先例、文本與解釋的枝節問題轉移權利爭點，屬結果面的理由，第肆部分展開；第二條見\mbox{ K05。}}
\keyentry{2}{K05}{1353}{}{而就民主價值而言，它在政治上欠缺正當性：它以少數幾位未經選舉、不受課責的法官之間的多數表決為準，在權利爭點的終局解決上，剝奪了一般公民的參政權，也把代表與政治平等這些向來受珍視的原則\mbox{撥到一旁。}}{And it is politically illegitimate, so far as democratic values are concerned: By privileging majority voting among a small number of unelected and unaccountable judges, it disenfranchises ordinary citizens and brushes aside cherished principles of representation and political equality in the final resolution of issues about rights.}{程序面理由在導論中的完整表述，第伍\mbox{部分展開。}}
\keyentry{3}{K88}{1353}{}{我反對違憲審查的論證不是無條件的，而是取決於現代自由民主體制的某些制度與\mbox{政治特徵。}}{My argument against judicial review is not unconditional but depends on certain institutional and political features of modern liberal democracies.}{七個部分的分工預告，此句交代第貳部分的作用；論證附條件的正式表述在\mbox{ K09。}}
\subsection*{壹、違憲審查的界定}
\keyentry{3}{K51}{1353}{}{本文談的是對法律的\mbox{違憲審查}}{This is an Essay about judicial review of legislation, not judicial review of executive action or administrative decisionmaking.}{本文的標的限於對法律的違憲審查；法院依法治原則審查行政行為，幾乎無人爭執，不在\mbox{本文範圍。}}
\keyentry{1}{K06}{1354}{一級・標的}{最重要的區別，在於我所稱的強式違憲審查與弱式違憲審查。我攻擊的標的是強式\mbox{違憲審查。}}{The most important difference is between what I shall call strong judicial review and weak judicial review. My target is strong judicial review.}{攻擊對象限於強式違憲審查。英國 1998 年人權法、紐西蘭權利法案所採的弱式審查不在此列\pg{1355}，討論時常把兩者\mbox{混為一談。}}
\keyentry{2}{K07}{1354}{}{在強式違憲審查的制度下，法院有權在個案中拒絕適用某部制定法（即使該法依其文義明顯適用於該案），或修改制定法的效果，使其適用合於個人權利（以該法本身未預\mbox{見的方式）。}}{In a system of strong judicial review, courts have the authority to decline to apply a statute in a particular case (even though the statute on its own terms plainly applies in that case) or to modify the effect of a statute to make its application conform with individual rights (in ways that the statute itself does not envisage).}{強式審查的定義：法院得於個案拒絕適用法律，或修\mbox{改其效果。}}
\keyentry{3}{K52}{1354}{}{而且，這種制度下的法院有權在法律上確立某部制定法或某項立法規定不予適用，於是因遵循先例原則與爭點效之故，它們拒絕適用的法律事實上\mbox{成為具文。}}{Moreover, courts in this system have the authority to establish as a matter of law that a given statute or legislative provision will not be applied, so that as a result of stare decisis and issue preclusion a law that they have refused to apply becomes in effect a dead letter.}{補充 K07 的強式審查定義：經遵循先例原則與爭點效，被拒絕適用的法律事實上\mbox{成為具文。}}
\keyentry{3}{K186}{1354}{}{還有一種更強的違憲審查形式，則授權法院逕行撤銷一項法律，使其自現行法中\mbox{完全除去。}}{A form of even stronger judicial review would empower the courts to actually strike a piece of legislation out of the statute-book altogether.}{強弱光譜的最強一端：法院得將法律自法典中除去；美國法院依原文\mbox{似無此權。}}
\keyentry{2}{K08}{1355}{}{相對地，在弱式違憲審查的制度下，法院得審查法律是否合於個人權利，但不得僅以適用該法將侵害權利為由，拒絕適用（或限\mbox{縮其適用）。}}{In a system of weak judicial review, by contrast, courts may scrutinize legislation for its conformity to individual rights but they may not decline to apply it (or moderate its application) simply because rights would otherwise be violated.}{弱式審查的定義，與 K07\mbox{ 對照讀。}}
\keyentry{3}{K89}{1356}{}{還有一種更弱的違憲審查形式，甚至不給法官這\mbox{樣的權限。}}{A form of even weaker judicial review would give judges not even that much authority.}{紐西蘭的形態：法院不得拒絕適用法律，只能尋求合乎權利的解釋，其不相容宣告對立法程序不生法律上的效果。與 K08 的弱式\mbox{審查並讀。}}
\keyentry{3}{K53}{1356}{}{但加拿大的法律（省或聯邦的法律）得以特定形式作成，使其免受這種審查——加拿大的議會得於法律中宣告，該法「儘管」有憲章所定的權利，\mbox{仍應施行。}}{But Canadian legislation (provincial or national) may be couched in a form that insulates it from this scrutiny—Canadian assemblies may legislate ``notwithstanding'' the rights in the Charter.}{加拿大議會得以「儘管條款」使法律免受審查，實務上甚少援用；作者因此把加拿大的安排算作強式\mbox{違憲審查。}}
\keyentry{3}{K149}{1356}{}{然而在實務上，「儘管條款」甚\mbox{少被援用。}}{In practice, however, the notwithstanding clause is rarely invoked.}{加拿大設有立法推翻的條款而仍歸強式，理由在此句：該條款\mbox{甚少援用。}}
\keyentry{3}{K54}{1358}{}{但有兩點要記住：以結構性限制為取向的進路，作為思考權利的依據可能不是特別合適；此外，許多針對以權利為取向之違憲審查所提的質疑，同樣可以對其他形式的憲法\mbox{審查提出。}}{But it is important to remember both that an approach oriented to structural constraints might not be particularly appropriate as a basis for thinking about rights, and the additional point that many of the challenges to rights-oriented judicial review can be posed to other forms of constitutional review as well.}{以結構性限制為取向的審查不宜移作思考權利的依據；對以權利為取向之審查的質疑，也及於聯邦主義等結構性的\mbox{憲法審查。}}
\keyentry{4}{K203}{1358}{四級・生動表述}{於是它改由大法官之間的表決來解決：有些人投給一種聯邦主義的構想（然後把它讀進憲法裡），其餘的人投給另一種，在法院裡票數較多的\mbox{一方勝出。}}{So it is settled instead by voting among Justices—some voting for one conception of federalism (which they then read into the Constitution), the others for another, and whichever side has the most votes on the Court prevails.}{括號裡「然後把它讀進憲法裡」揭穿解釋的次序：先投票選定構想，再說它在憲法裡；可預告第伍部分 K35 的\mbox{「數人頭」。}}
\keyentry{3}{K55}{1359}{}{對某些為違憲審查所作的辯護而言，違憲審查在事後行使、扎\mbox{根於個案}}{For some defenses of judicial review, the a posteriori character of its exercise—its rootedness in particular cases}{本文聚焦於事後、扎根於個案的審查，因為部分辯護以此一性質為據；第肆部分回應「法院面對具體個案」的主張，見\mbox{ K26。}}
\keyentry{3}{K56}{1359}{}{一般審級體系裡的法官就權利進行推理的能力被高估了，因為審判日常的種種紀律，多半把他們的注意力從道德論證\mbox{本身引開。}}{The ability of judges in the regular hierarchy of courts to reason about rights is exaggerated when so much of the ordinary discipline of judging distracts their attention from direct consideration of moral arguments.}{審判的日常紀律使法官的注意力離開道德論證；第肆部分論文本形式主義（K28\mbox{）時展開。}}
\keyentry{4}{K204}{1359}{四級・生動表述}{專門的憲法法院或許能做得好一些，不過經驗顯示，它也可能過度專注於自身學理與先例的發展，從而使它審酌以權利為本的推理時，隔著一層扭\mbox{曲的濾鏡。}}{Perhaps a specialist constitutional court can do better, though experience suggests that it too may become preoccupied with the development of its own doctrines and precedents in a way that imposes a distorting filter on the rights-based reasoning it considers.}{「扭曲的濾鏡」比喻學理與先例擋在法院與道德論證之間；導讀時可連到第肆部分法院說\mbox{理的批評。}}
\subsection*{貳、四項假設}
\keyentry{3}{K150}{1360}{}{其中某些假設在有些讀者看來可能是在乞題（question-begging），但我在此\mbox{無意取巧。}}{Certain of these assumptions may strike some readers as question-begging, but I am not trying any sort of subterfuge here.}{作者預先回應「四項假設是乞題」的質疑，下文接著說非核心情形留待\mbox{第柒部分。}}
\keyentry{1}{K09}{1360}{一級・條件}{實際上，我的主張是：反對違憲審查的論證附有條件；任何一項條件不成立，論證便可\mbox{能不成立。}}{In effect, my contention will be that the argument against judicial review is conditional; if any of the conditions fail, the argument may not hold.}{論證以四項假設為前提，任一假設不成立，結論即可能不成立。第柒部分的非核心情形以此\mbox{句為起點。}}
\keyentry{4}{K205}{1360}{四級・生動表述}{我完全預料到，有些讀者會快速翻檢我的假設，看哪幾項不適用於他們所理解的社會，比方說美國或英國社會，找到了便把核心論證整個擱置不顧。對他們而言，要緊的是違憲審查得到辯護、對它的質疑能被擊退；至於怎麼做到，他們並不\mbox{特別在意。}}{I fully expect that some readers will comb quickly through my assumptions to find some that do not apply, say, to American or British society as they understand it, leading them to ignore the core argument altogether. What matters to them is that judicial review be defended and challenges to it seen off; they don't particularly care how.}{作者預先描繪「只問結論」的讀者如何翻檢假設、找到一項就丟下論證；導讀時用來提醒課堂不要這樣讀\mbox{四項假設。}}
\keyentry{2}{K90}{1360}{}{我們要設想一個社會，它有：（1）運作良好的民主制度，包括一個依成年人普選產生的代議立法機關；（2）一套同樣運作良好的司法制度，以非代議的方式設置，用以審理個別訴訟、解決爭議並維護法治；（3）該社會多數成員與多數公職人員信奉個人與少數者享有權利此一理念；以及（4）在信奉權利理念的社會成員之間，關於權利（亦即關於對權利的信奉究竟意味著什麼、又有何蘊涵）持續、實質而出於善\mbox{意的分歧。}}{We are to imagine a society with (1) democratic institutions in reasonably good working order, including a representative legislature elected on the basis of universal adult suffrage; (2) a set of judicial institutions, again in reasonably good order, set up on a nonrepresentative basis to hear individual lawsuits, settle disputes, and uphold the rule of law; (3) a commitment on the part of most members of the society and most of its officials to the idea of individual and minority rights; and (4) persisting, substantial, and good faith disagreement about rights (i.e., about what the commitment to rights actually amounts to and what its implications are) among the members of the society who are committed to the idea of rights.}{四項假設的總述。K09 說論證以這些假設為條件，第貳部分其餘各段逐項展開，第柒部分討論假設不成\mbox{立的情形。}}
\keyentry{1}{K10}{1360}{一級・正面主張}{我將主張，在這些假設之下，該社會應當用它的立法制度來解決其成員之間關於權\mbox{利的分歧。}}{I shall argue that, relative to these assumptions, the society in question ought to settle the disagreements about rights that its members have using its legislative institutions.}{四項假設之下，權利分歧交由立法制度解決。其後各部比較兩種決定程序，所支持的都是\mbox{這個主張。}}
\keyentry{2}{K11}{1361}{}{我假設這個立法機關是一個大型的審議機關，慣於處理困難的議題，包括正義與社會政策的\mbox{重大議題。}}{I assume that this legislature is a large deliberative body, accustomed to dealing with difficult issues, including important issues of justice and social policy.}{第一項假設所設想的立法機關：大型、慣於審議正義與社會政策的\mbox{重大議題。}}
\keyentry{2}{K151}{1361}{}{整體而言，我假設民主制度運作良好。它們未必完美，關於如何改進，大概也有持\mbox{續的討論。}}{In general, I am assuming that the democratic institutions are in reasonably good order. They may not be perfect and there are probably ongoing debates as to how they might be improved.}{第一項假設的限度：只要求運作良好，不要求完美。聽者質疑立法程序被理想化時，以\mbox{此句回答。}}
\keyentry{2}{K91}{1361}{}{我假設這些討論受一種民主文化的浸潤，這種文化重視負責任的審議與\mbox{政治平等。}}{I assume these debates are informed by a culture of democracy, valuing responsible deliberation and political equality.}{第一項假設的內容。同段接著強調政治平等，第伍部分程序面的論證以政治平\mbox{等為依據。}}
\keyentry{3}{K12}{1362}{}{一旦設定了這第一項假設，還能有什麼理由要另設一套非民選的程序，來審查、有時推翻立法機關已完\mbox{成的工作？}}{Once we have posited this first assumption, what reason can there be for wanting to set up a nonelective process to review and sometimes override the work that the legislature has done?}{第一項假設設定後的提問。後文各部都在回答：何以要另設非民選的程序審查、\mbox{推翻立法。}}
\keyentry{3}{K152}{1362}{}{另一方面，我不想以這個起始的假設乞題。我隨即以另一項假設與之平衡：我們所設定的這個社會同樣有運作良好的法院，能把法院擅長的事做得相當好。這將是第\mbox{二項假設。}}{On the other hand, I do not want to beg any questions with this initial assumption. I shall balance it immediately with the assumption that the society we are postulating also has courts in good working order—this will be the second assumption—doing reasonably well what courts are good at doing.}{第一、二項假設對稱設定，兩個機關都以運作良好為前提，比較才不\mbox{偏向一方。}}
\keyentry{3}{K57}{1362}{}{我關於立法與選舉安排運作良好所說的一切，涉及的是程序面的價值，而非結果\mbox{面的價值。}}{All that I have said about the legislative and electoral arrangements being in good working order goes to process values rather than outcome values.}{「運作良好」只就程序而言，不假設立法機關所制定的法律在內容上\mbox{大體正義。}}
\keyentry{2}{K58}{1363}{}{我假設，與前一小節所指的機關不同，法院大體上不是民選或代\mbox{議的機關。}}{I assume that, unlike the institutions referred to in the previous Section, the courts are mostly not elective or representative institutions.}{第二項假設下的法院大體上非民選、非代議；第伍部分比較立法選舉與司法任命的正當性（K34）以\mbox{此為前提。}}
\keyentry{4}{K206}{1363}{四級・生動表述}{許多為違憲審查辯護的人視此為一大優點，因為這表示法院得以就原則問題進行審議，不受民意壓力的干擾，也不為公眾的怒氣所動。不過有時，當有人認為有必要駁斥從民主出發而反對違憲審查的論據時，這個制度的辯護者會自豪地指出法官由選舉\mbox{產生的州。}}{Many defenders of judicial review regard this as a huge advantage, because it means courts can deliberate on issues of principle undistracted by popular pressures and invulnerable to public anger. Sometimes, however, when it is thought necessary to rebut the democratic case against judicial review, defenders of the practice will point proudly to states where judges are elected.}{同一批辯護者先誇法官不受民意干擾，被質疑民主正當性時又「自豪地」指出民選法官；兩句並讀\mbox{即見反諷。}}
\keyentry{3}{K13}{1363}{}{法院能夠審查法律；問題在於它們是否應當如此，以及若應當，其判斷是否應為終局並拘束政府的\mbox{代議部門。}}{They could review legislation; the question is whether they should, and if so, whether their determinations should be final and binding on the representative branches of government.}{爭點限定在法院應否審查立法、其判斷應否終局拘束代表部門；法院有無審查的能力不在\mbox{爭執之列。}}
\keyentry{2}{K14}{1364}{}{也就是說，我假設關於權利分歧的第四項假設同樣適用於法官，而這也會影響他們行使違憲審查權的方式（如果他們有這項\mbox{權力的話）。}}{That is, I assume they are subject to my fourth assumption about rights-disagreement, and that this too affects how they exercise powers of judicial review (if they have such powers).}{第四項假設同樣適用於法官，法官之間也有\mbox{權利分歧。}}
\keyentry{2}{K59}{1364}{}{我假設我們所設想的這個社會，其多數成員堅定信奉個人與少數者享有權利\mbox{此一理念。}}{I assume that there is a strong commitment on the part of most members of the society we are contemplating to the idea of individual and minority rights.}{第三項假設的表述：社會多數成員堅定信奉個人與少數者的權利。K15、K1\mbox{6 承此。}}
\keyentry{2}{K92}{1365}{}{我假設對權利的信奉不只是口頭上的表態，社會成員認真\mbox{對待權利：}}{I assume that the commitment to rights is not just lip service and that the members of the society take rights seriously:}{第三項假設的內容：這種信奉須是真實的，社會成員持續辯論權利並留意一切社會決定中的權利問題。K5\mbox{9 承此。}}
\keyentry{3}{K93}{1365}{}{但我假設，對個人與少數者的權利的一般尊重，是該社會廣泛共識中嚴肅的一部分，是最通行的那一套政治意見的一部分，也確實是官方意識形態\mbox{的一部分。}}{But I assume that general respect for individual and minority rights is a serious part of a broad consensus in the society, part of the most prevalent body of political opinion, and certainly part of the official ideology.}{第三項假設的補充：社會裡懷疑權利的人\mbox{屬於例外。}}
\keyentry{3}{K15}{1365}{}{在這個階段我只假設，該社會已制定一部權利法案，以體現它對權\mbox{利的信奉。}}{All I assume at this stage is that a Bill of Rights has been enacted to embody the society's commitment to rights.}{第三項假設只要求社會已制定權利法案，以體現對權\mbox{利的信奉。}}
\keyentry{4}{K207}{1366}{四級・生動表述}{一方面，我似乎在做不利於自己的論證：假設有一部權利法案，等於事先做了手腳，讓局面有利於違憲審查。另一方面，這看來又像是要耍\mbox{什麼花招。}}{On the one hand, I appear to be arguing against interest, stacking the deck in favor of judicial review by assuming a Bill of Rights. On the other hand, it may seem that something sneaky is in the offing.}{牌戲作弊的比喻（stack the deck）與「耍花招」，作者自揭兩面都可能被懷疑；第伍部分 K148 再用\mbox{同一比喻。}}
\keyentry{2}{K16}{1366}{}{我聽膩了反對違憲審查的人被貶為權利懷疑論者。最好的回應，是把反對違憲審查的論據建立在對權利堅定而普遍的\mbox{信奉之上。}}{I am tired of hearing opponents of judicial review denigrated as being rights-skeptics. The best response is to erect the case against judicial review on the ground of a strong and pervasive commitment to rights.}{回應「反對違憲審查者是權利懷疑論者」的指摘：本文的論據以社會對權利的堅定信\mbox{奉為前提。}}
\keyentry{3}{K153}{1366}{}{這第三項假設把對權利的信奉薄弱而脆弱的社會界定為非\mbox{核心情形。}}{This third assumption defines as non-core cases societies in which the commitment to rights is tenuous and fragile.}{第三項假設不成立的社會歸入非核心情形，四項假設表「不成立時」一\mbox{欄由此填。}}
\keyentry{2}{K94}{1366}{}{這類論證饒富興味，但它們並未觸及美國、英國、加拿大這類國家為違憲審查所提論\mbox{據的核心。}}{Such arguments are interesting, but they do not go to the heart of the case that is made for judicial review in countries like the United States, Britain, or Canada.}{第三項假設把對權利的信奉脆弱的社會列為非核心情形。辯護者說違憲審查可以支撐對權利的信奉，作者答以美、英、加三國的論據是把社會對權利既有的信\mbox{奉制度化。}}
\keyentry{2}{K95}{1366}{}{我最後一項、也是關鍵的一項假設是：在現代自由社會中，我們看到一切重大政治議題都有普遍的分歧，關於權利的共識也不能免於\mbox{這種分歧。}}{My final and crucial assumption is that the consensus about rights is not exempt from the incidence of general disagreement about all major political issues, which we find in modern liberal societies.}{第四項假設的內容。同段把分歧分為哲學層面、立法提案、困難個案\mbox{三個層次。}}
\keyentry{2}{K96}{1367}{}{我假設權利分歧大多不是狹義法條主義意義上的\mbox{解釋問題。}}{I assume that the rights-disagreements are mostly not issues of interpretation in a narrow legalistic sense.}{第四項假設的分歧屬於「分水嶺式」的政治哲學議題，墮胎、優惠性差別待遇等例子都屬此類。K1\mbox{7 承此。}}
\keyentry{2}{K17}{1367}{}{即使分水嶺式的案件仍有爭議，對權利的信奉仍可以是全心\mbox{而真誠的。}}{A commitment to rights can be wholehearted and sincere even while watershed cases remain controversial.}{對權利的信奉與分水嶺案件上的分歧可以並存，第四項假\mbox{設的核心。}}
\keyentry{4}{K208}{1368}{四級・生動表述}{無疑有無賴（對權利毫不在意）或道德文盲（誤解權利的效力與重要性）以虛偽或無知的態度持守並辯護\mbox{某些立場。}}{No doubt some positions are held and defended disingenuously or ignorantly by scoundrels (who care nothing for rights) or moral illiterates (who misunderstand their force and importance).}{「無賴」與「道德文盲」兩個辛辣的稱呼，括號各附定義；導讀時用來說明作者承認例外，再把它們排出\mbox{核心情形。}}
\keyentry{2}{K97}{1368}{}{但我假設多數情形下，人們的分歧是理性而出\mbox{於善意的。}}{But I assume that in most cases disagreement is pursued reasonably and in good faith.}{第四項假設中「出於善意」的部分：有分歧不表示一方不認真對待權利，作者援用 Rawls 的「判斷的負擔」並及於正\mbox{當的議題。}}
\keyentry{3}{K154}{1368}{}{換言之，我所假設的近於 John Rawls 所稱的「判斷的負擔」，只是同時適用於正當的議題與善的議題；正當的議題，正是 Rawls 遲疑於適用該\mbox{學說之處。}}{In other words, I assume something like John Rawls's ``burdens of judgment,'' but applied (where Rawls hesitated to apply the doctrine) to issues of the right as well as issues of the good.}{善意分歧的理論依據取自 Rawls，並擴及\mbox{權利議題。}}
\keyentry{3}{K18}{1368}{}{分歧的假設與道德相對\mbox{主義無關。}}{The assumption of disagreement has nothing to do with moral relativism.}{權利分歧的假設與道德相對\mbox{主義無涉。}}
\keyentry{2}{K19}{1368}{}{如果有一部權利法案，我假設它與分歧中的爭點有關，但不解決\mbox{這些爭點。}}{If there is a Bill of Rights, I assume that it bears on, but does not resolve, the issues at stake in the disagreements.}{權利法案與分歧的爭點有關而不解決之。第肆部分 B 節論法院依附文本時，\mbox{取用此句。}}
\keyentry{4}{K209}{1369}{四級・生動表述}{事實上，律師將大有文章可做。每項分歧的每一方都會主張，只要把那些文本讀得夠寬（或夠窄），自己的立場就可以讀進權利法案那些空泛\mbox{的承諾裡。}}{In fact, lawyers will have a field day. Each side to each of the disagreements will claim that its position can be read into the bland commitments of the Bill of Rights if only those texts are read generously (or narrowly) enough.}{「律師的大好日子」與括號裡的「或夠窄」，寫兩造都能把立場讀進同一條空泛條文；導讀時接\mbox{ K98。}}
\keyentry{4}{K210}{1369}{四級・生動表述}{雙方都不會願意公開承認一件我認為此刻已顯而易見的事：認真對待權利的人之間必有真實而合理的分歧；權利法案空泛的措辭，設計上不過是要把這些分歧巧妙迴避過去，撐到法案\mbox{制定為止。}}{Neither will be prepared to acknowledge publicly that which I am assuming now will be obvious: that the bland rhetoric of the Bill of Rights was designed simply to finesse the real and reasonable disagreements that are inevitable among people who take rights seriously for long enough to see such a Bill enacted.}{權利法案的措辭被說成刻意含糊、只求撐到通過的妥協文字；可與第肆部分 K2\mbox{7 對讀。}}
\keyentry{3}{K98}{1369}{}{違憲審查不會鼓勵我們直接面對這些分歧，反而很可能使這些分歧被當成那些空泛表述的解釋問\mbox{題來處理。}}{Instead of encouraging us to confront these disagreements directly, judicial review is likely to lead to their being framed as questions of interpretation of those bland formulations.}{承 K19：權利法案的措辭空泛，違憲審查使分歧轉成文本解釋問題。第肆部分 K27、K2\mbox{8 展開。}}
\subsection*{參、論證的形式}
\keyentry{4}{K187}{1370}{四級・生動表述}{有各式各樣的問題，並不需要全社會得出一個定論，例如聖餐變體、《哈姆雷特》的意義、純粹沉思生活的價值；這是幸運的，因為在這些領域幾乎沒有達成一致的指望。不幸的是，在我們確實需要解決的權利問題上，同樣幾乎沒有達成一\mbox{致的指望。}}{Now, there are all sorts of issues on which we do not need society-wide settlement—transubstantiation, the meaning of Hamlet, the value of a purely contemplative life—and that is fortunate, because there is little prospect of agreement in these areas. Unfortunately, on issues of rights, for which we do need settlement, there is little prospect of agreement either.}{以聖餐變體、《哈姆雷特》對照權利問題：前者不需全社會的定論，權利問題需要，卻同樣\mbox{難有共識。}}
\keyentry{3}{K60}{1370}{}{需要作成決定，並不會使分歧不復存在；毋寧是說，共同行動的基礎必須在激烈的分歧之中\mbox{打造出來。}}{The need for settlement does not make the fact of disagreement evaporate; rather, it means that a common basis for action has to be forged in the heat of our disagreements.}{作成決定的必要並不消除分歧，共同行動的基礎須在分歧之中形成；K20 對決定程序的形式要求以\mbox{此為前提。}}
\keyentry{3}{K99}{1370}{}{在現實世界裡，必須作成決定的時刻出現在立\mbox{法的場域。}}{In the real world, the need for settlement confronts us in the legislative arena.}{權利問題通常在法律適用時才浮現，必須決定的是該法律應否依其文義適用。K61 的提問以\mbox{此為背景。}}
\keyentry{2}{K100}{1370}{}{但這只是支持弱式違憲審查的論證，並不支持這個制度的強式形態；在那種形態下，已被指認的抽象權利問題依法院認為適當的\mbox{方式解決。}}{But this is an argument for weak judicial review only, not for a strong form of the practice in which the abstract question of right that has been identified is settled in the way that a court deems appropriate.}{作者對「議員未必看得出權利問題」這項辯護的回應：該理由只支持英國式的宣告或檢察總長的審查，支持不了\mbox{強式審查。}}
\keyentry{3}{K61}{1371}{}{我們面對的問題是：立法機關的這項決議應否具有決定性，還是有理由讓它由司法機關事後複核，或許還\mbox{予以推翻。}}{The question we face is whether that resolution of the legislature should be dispositive or whether there is reason to have it second-guessed and perhaps overruled by the judiciary.}{立法機關已就權利問題審議並表決，爭點在於其決議應否由司法機關複核、推翻。第參部分以下比較兩類理由，回答的都\mbox{是這一問。}}
\keyentry{2}{K20}{1371}{}{我們必須建立一套決定程序，其運作要平息、而非重新點燃那些當初使決定程序成為必\mbox{要的爭議。}}{We must set up a decision-procedure whose operation will settle, not reignite, the controversies whose existence called for a decision-procedure in the first place.}{決定程序的形式要求，引 Hobbes 為據\pg{1371}。「立法有效，除非侵害權利」因此不能充當決定規則：社會成員爭執的正是某法律是否\mbox{侵害權利。}}
\keyentry{2}{K21}{1371}{}{這表示，即使我們所設想的社會成員對權利有分歧，他們仍需就用以解決分歧的決定程序共享一套正\mbox{當性理論。}}{This means that even though the members of the society we are imagining disagree about rights, they need to share a theory of legitimacy for the decision-procedure that is to settle their disagreements.}{決定程序的第二項形式要求：成員須共享一套關於程序正當\mbox{性的理論。}}
\keyentry{2}{K101}{1372}{}{另一方面，有一些與正當性有關的重要理由——例如公平、發言權、參與——是因分歧而生的，若不處理決定程序的問題就\mbox{不會出現。}}{For another thing, there are important reasons relating to legitimacy—e.g., fairness, voice, participation—that arise because of disagreement and do not arise apart from our addressing the question of decision-procedures.}{對「既然程序也有分歧，何不直接就實質採取立場」的回應：不論有無分歧都須處理正當性問題，公平、發言權、參與等理由只在處理決定程序\mbox{時才出現。}}
\keyentry{2}{K102}{1372}{}{沒有一套決定程序會是完美的。不論它是不受審查的立法程序，還是違憲審查的程序，它有時都會作成錯誤的決定，背棄權利而非\mbox{維護權利。}}{No decision-procedure will be perfect. Whether it is a process of unreviewable legislation or whether it is a process of judicial review, it will sometimes come to the wrong decision, betraying rights rather than upholding them.}{兩種程序都會出錯，比較的是相對的優劣。同段把 Wollheim 的弔詭改定為政治理論的一般弔詭，第陸部分的 P49 \mbox{回到此點。}}
\keyentry{3}{K188}{1372}{}{但 Wollheim 把這個弔詭歸給民主是錯的。凡是政治理論，只要在說明應當做什麼之外，另行說明人們對此有分歧時應如何作成決定，都受這個弔詭影響；它是政治理論的\mbox{一般弔詭。}}{But Wollheim was wrong to ascribe this paradox to democracy. It is a general paradox in the theory of politics affecting any political theory that complements its account of what ought to be done with an account of how decisions ought to be made when there is disagreement about what ought to be done.}{Wollheim 弔詭被作者改寫為一切政治理論的弔詭，違憲審查同樣\mbox{受其影響。}}
\keyentry{2}{K62}{1372}{}{程序面的理由，是堅持應由某個人作成或參與作成某項決定的理由，其成立與何種結果才算適當的\mbox{考量無涉。}}{Process-related reasons are reasons for insisting that some person make, or participate in making, a given decision that stand independently of considerations about the appropriate outcome.}{程序面理由的定義：要求由某人作成或參與決定，與結果是否適當無涉；政治上的典型是政治平等與投票權。第伍\mbox{部分展開。}}
\keyentry{3}{K155}{1372}{}{在個人生活中，我們有時說，父母有權決定子女是否應因某次過犯受管教：這不是街上的路人或公車上另一位乘客可以\mbox{決定的事。}}{In personal life, we sometimes say that a parent has the right to make the decision as to whether her child should be disciplined for a given infraction: It is not for a passer-by on the street or another passenger on the bus to make that decision.}{程序面理由的設例：決定權歸屬與決定\mbox{好壞無關。}}
\keyentry{2}{K63}{1373}{}{相對地，結果面的理由，是要求決定程序的設計能確保適當結果（亦即好的、正義的或正確的決定\mbox{）的理由。}}{Outcome-related reasons, by contrast, are reasons for designing the decision-procedure in a way that will ensure the appropriate outcome (i.e., a good, just, or right decision).}{結果面理由的定義，與 K62 並讀。在原則事項上答錯，權利即受侵害，結果面理由因此須認真看待。第肆\mbox{部分展開。}}
\keyentry{2}{K103}{1373}{}{既然有分歧，關鍵就在於建立一套程序，其所產生的解決，雙方都能承認是以正當方\mbox{式作成的。}}{Given the disagreement, the whole point here is to set up a procedure for generating settlements in a way that can be recognized as legitimate on both sides.}{結果面理由的限制：不得以偏向某一爭議結果為由選擇程序，否則對立的一方不會接受。承\mbox{ K20。}}
\keyentry{2}{K156}{1373}{}{我們不必以乞題的方式說，應當選擇最可能產出某一組特定爭議性權利的政治程序，而可以改說我們應當選擇最可能觸及權利之真相的政治程序，不論那個真相最\mbox{終是什麼。}}{Instead of saying (in a question-begging way) that we should choose those political procedures that are most likely to yield a particular controversial set of rights, we might say instead that we should choose political procedures that are most likely to get at the truth about rights, whatever that truth turns out to be.}{結果面理由可接受的版本：以觸及權利真相為目標，不預設特定結果。第肆部分的比較依此\mbox{版本進行。}}
\keyentry{3}{K104}{1373}{}{許多工具主義的論證並不建立在對任何特定權利之內容\mbox{的知識上。}}{Many instrumentalist arguments are not based on knowledge of the content of any particular rights.}{本段全為引文（Kavanagh）：結果面論證可以只憑一般性的制度考量提出，不必先確定權\mbox{利的內容。}}
\keyentry{3}{K64}{1374}{}{Raz 沒有充分注意到這一點：儘管結果面的理由在關於權利的這個決定領域裡非常重要，其他種類的理由也\mbox{可能重要。}}{Raz has paid insufficient attention to the point that although outcome-related reasons are very important in this area of decisionmaking about rights, reasons of other kinds may be important too.}{回應 Raz「只有結果面理由值得考慮」之說，同段以自決\mbox{原則為例。}}
\keyentry{3}{K105}{1374}{}{但他們的責任仍是以權利為本的責任——在這方面有一項\mbox{注意義務。}}{But their responsibility is still a rights-based responsibility—there is a duty to take care in this regard.}{結果面理由的規範性質：避免權利侵害帶有義務論上的急迫性，設計者只負間接責任，仍屬以權利為本的注意義務。與 K65 程序面理由的\mbox{性質並讀。}}
\keyentry{2}{K65}{1375}{}{程序面的理由往往同樣是具有義務論急\mbox{迫性的事。}}{Process-related reasons are often matters of deontological urgency also.}{程序面理由同具義務論上的急迫性；作者據此批評 Dworkin 把參與的理由說成參與帶\mbox{來的後果。}}
\keyentry{4}{K189}{1375}{四級・生動表述}{這一切無疑重要。但這聽來像一位校長在列舉：設學生會、讓學生對校務有發言機會，可能帶來\mbox{哪些好處。}}{All this is no doubt important. But it has the flavor of a headmaster noting the advantages that may accrue from giving his pupils a say in educational affairs through a school council.}{校長與學生會之喻：Dworkin 把參與當成帶來好處的安排；下一句 K157 說這低估了參\mbox{與的權利。}}
\keyentry{3}{K157}{1375}{}{Dworkin 的說明嚴重低估了參與之權利這個概念，低估了「在社會的決定上，人人應受平等對待」這項要求，也低估了當有人憤慨地質問「這個決定影響我，我又必須服從它，他們怎麼敢把我的發言排除在外——剝奪我的參政權？」時，那種事關原\mbox{則的感受。}}{Dworkin's account radically underestimates the notion of a right to participate, the imperative that one be treated as an equal so far as a society's decisionmaking is concerned, the sense of principle that is at stake when someone asks indignantly, ``How dare they exclude my say—disenfranchise me—from this decision, which affects me and to which I am subject?''}{回應 Dworkin 把參與說成「參與上的後果」：參與是權利，不只是\mbox{帶來好處。}}
\keyentry{4}{K211}{1375}{四級・生動表述}{我們面對一個熟悉的問題：要同時讓兩個變數的值都最大，就像要人以最低的價格買\mbox{最快的車。}}{We face the familiar problem of trying to maximize the value of two variables, like asking someone to buy the fastest car at the lowest price.}{「最便宜又最快的車」比喻程序面與結果面兩種理由難以同時極大化；導讀時用來引出作\mbox{者的解法。}}
\keyentry{4}{K212}{1375}{四級・生動表述}{我想我可以快刀斬亂麻，了結\mbox{這個難題。}}{I think I can cut through this Gordian knot.}{Gordian knot 的典故（亞歷山大一劍斬斷戈爾迪之結），譯文改用「快刀斬亂麻」；接著 K22、K23 即是作\mbox{者的解法。}}
\keyentry{1}{K22}{1375}{一級・結果面}{我在第肆部分要論證的是，結果面的理由至多不\mbox{具決定性。}}{What I will argue, in Part IV, is that the outcome-related reasons are at best inconclusive.}{第肆部分的結論先在此宣告。法院被歸予的三項優點（面對個案、依據文本、說明理由），作者在第肆部分\mbox{逐一回應。}}
\keyentry{1}{K23}{1375}{一級・程序面}{程序面的理由則相當一面倒。它們主要的作用是使違憲審查失去信用，卻不傷及\mbox{立法決定。}}{The process-related reasons, however, are quite one-sided. They operate mainly to discredit judicial review while leaving legislative decisionmaking unscathed.}{第伍部分的結論。與 K22 合讀，即篇名所稱 the core of the case\pg{1375--1376}。}
\subsection*{肆、結果面的理由}
\keyentry{3}{K106}{1376}{}{但如 Lon Fuller 所指出，由此並不能推出：法院因此就是處理此處所考慮之權利主張（取其較不具技術性的意義）的\mbox{適當場所。}}{But as Lon Fuller observed, it does not follow that courts are therefore the appropriate forum for dealing with claims of right in the less technical sense under consideration here.}{回應 Raz 以執行能力決定機關的說法：法院只擅長兩造對立的問題，多面向的權利問題不宜交付司法。作者在\mbox{此不展開。}}
\keyentry{3}{K24}{1376}{}{下文將會說明，結果面的理由同時指向\mbox{兩個方向。}}{Outcome-related reasons, as we shall see, cut in both directions.}{第肆部分的總論：結果面的理由，支持與反對違憲審查的兩方\mbox{都可援引。}}
\keyentry{3}{K107}{1377}{}{因此我們可能偏好這樣的程序：受託作成決定的人，不受或不直接受其自身決\mbox{定的影響。}}{We may therefore prefer a procedure in which those charged with a decision are not affected, or not directly affected, by their own decision.}{本段全為 Raz 的引文：立法者受派系利益驅動，宜由不受決定影響的人決定。下\mbox{一段回應。}}
\keyentry{2}{K108}{1377}{}{同樣的派系壓力也常常可以解釋司法對權\mbox{利的忽視。}}{The same sectarian pressures often explain judicial neglect of rights as well.}{對 Raz 的回應：派系偏見若普遍存在即屬非核心情形；即使照單全收，Korematsu 等案與 Tribe 對 9/11 之後的觀察顯示法院\mbox{同受其害。}}
\keyentry{4}{K213}{1377}{四級・生動表述}{在美國，我們已從 Korematsu、Schenck、Dred Scott 與 Prigg 這些性質各異的案件中看\mbox{到這一點。}}{We have seen this in the United States in cases as diverse as Korematsu, Schenck, Dred Scott, and Prigg.}{連舉四件惡名昭彰的判決（日裔拘禁、反戰言論、奴隸、逃奴），以實例反駁法院較能抵擋\mbox{派系壓力。}}
\keyentry{4}{K190}{1377}{四級・生動表述}{較晚近，Laurence Tribe（通常是違憲審查的堅定辯護者）指出，在 9/11 之後籠罩美國的恐慌中，「對於憂心民權與自由的人來說，在危機時期把太多希望寄託在司法機關上，會是可\mbox{怕的錯誤。」}}{More recently, Laurence Tribe (usually a stalwart defender of judicial review) observed that in the panic that afflicted America after 9/11, ``it would be a terrible mistake for those who worry about civil rights and liberties to pin too much hope on the judiciary in times of crisis.''}{連違憲審查的支持者 Tribe 也說危機時期不宜寄望法院，用以佐證 \mbox{K108。}}
\keyentry{4}{K214}{1378}{四級・生動表述}{例如，它們將兒童排除於選舉權之外，儘管兒童也受到所討論的決\mbox{定的影響。}}{They exclude children from voting, for example, even though children are affected by the decisions under consideration.}{以排除兒童投票這個人人接受的例子，說明民主程序本來就為結\mbox{果而設計。}}
\keyentry{3}{K25}{1378}{}{我們在參與這一面所看到的，不是 Rawls 所稱的純粹的程序正義的主張，而是近似不完全的\mbox{程序正義。}}{What we see, on the participatory side, is not what Rawls called a claim of pure procedural justice, but something like imperfect procedural justice.}{參與權在結果面的意義，借 Rawls 不完全程序正義的\mbox{概念說明。}}
\keyentry{3}{K109}{1379}{}{人們勉力把結果面的理由與司法機關聯繫起來，而在這麼做的時候，對於司法決定如何作成，他們常兜售一幅相當不切實\mbox{際的圖像。}}{People strain to associate outcome-related reasons with the judiciary, and in so doing they often peddle a quite unrealistic picture of what judicial decisionmaking is like.}{作者對結果面辯護的總評，並說明反對者有時刻意以樂觀的立法機關圖像對應辯護者對法院的\mbox{樂觀圖像。}}
\keyentry{2}{K158}{1379}{}{反對違憲審查的人常被指責對立法機關採取一種天真樂觀的看法。但有時我們是故意的：法院的辯護者既然不肯如實說明法院裡發生的事，我們就以一幅樂觀的圖像對\mbox{應另一幅。}}{Opponents of judicial review are often accused of adopting a naively optimistic view of legislatures. But sometimes we do this deliberately, matching one optimistic picture with another in the face of the refusal of the defenders of courts to give a realistic account of what happens there.}{回應「把立法機關理想化」的指責：對立法機關的樂觀描述，是用來對應辯護者對法院的\mbox{樂觀描述。}}
\keyentry{3}{K66}{1379}{}{然而在這三點上，我都將論稱：法院處理權利的方式有重大的結果面缺陷，而立法機關這一方則有重大的結\mbox{果面優點。}}{On all three counts, however, I shall argue that there are important outcome-related defects in the way courts approach rights and important outcome-related advantages on the side of legislatures.}{法院被歸予的三項優點（面對個案、依據文本、說明理由），作者逐一回應的總綱；K26、K27、K29 \mbox{依序承之。}}
\keyentry{3}{K67}{1379}{}{有時有人論稱，司法對權利的推理（相對於立法對權利的推理）之美妙處在於，權利問題是以有血有肉的個別情境的形式呈現到\mbox{法官面前。}}{People sometimes argue that the wonderful thing about judicial reasoning on rights (as opposed to legislative reasoning on rights) is that issues of rights present themselves to judges in the form of flesh-and-blood individual situations.}{法院第一項優點的表述，同段引 Michael Moore 為例；K26\mbox{ 回應之。}}
\keyentry{4}{K215}{1379}{四級・生動表述}{如 Michael Moore 所陳述的：「法官在⋯⋯道德洞見上比立法機關處於較有利的位置，因為法官每天〔原文如此〕都面對道德思想實驗，其細節與具體的個人牽涉，正是道德洞見\mbox{所需要的。」}}{As Michael Moore puts the point, ``judges are better positioned for … moral insight than are legislatures because judges have moral thought experiments presented to them everyday [sic] with the kind of detail and concrete personal involvement needed for moral insight.''}{法官每天都在做「道德思想實驗」的說法，作者照錄並在 everyday 的拼法後加 [sic]，暗示不以為然；下一段 K26 即\mbox{斥為神話。}}
\keyentry{2}{K26}{1379--1380}{}{但這大體上是神話。違憲審查的爭論所談的，大多是高階的上訴審級；案件到達那裡時，原本有血有肉的權利人幾乎已不留痕跡，而所謂的論辯（如果算得上論辯的話）是繞著爭執中權利的抽象問\mbox{題進行的。}}{But this is mostly a myth. By the time cases reach the high appellate levels we are mostly talking about in our disputes about judicial review, almost all trace of the original flesh-and-blood right-holders has vanished, and argument such as it is revolves around the abstract issue of the right in dispute.}{回應法院的第一項優點（面對具體個案）：案件到了上訴審，只剩系爭權利的\mbox{抽象問題。}}
\keyentry{3}{K159}{1380}{}{的確，近來有一種趨向，是以惡名昭彰的個別案件為基礎發動立法，Megan's Law \mbox{即為一例。}}{Indeed, there is a tendency these days to initiate legislation on the basis of notorious individual cases—Megan's Law, for example.}{回應「法院以個案為取向」：立法程序同樣考量個案，舉 Megan's La\mbox{w 為例。}}
\keyentry{4}{K216}{1380}{四級・生動表述}{有人說，困難案件造就惡法（hard cases make bad\mbox{ law）}}{Hard cases make bad law, it is sometimes said.}{英美法諺「困難案件造就惡法」，作者借來翻轉：真要處理個案，立法機關反而較合適（\mbox{K110）。}}
\keyentry{2}{K110}{1380}{}{在這句話成立的範圍內，一般性的權利問題影響數百萬人，影響的方式又各不相同；我認為，要評估個別案件對這類問題的意義，立法機關處於遠較有\mbox{利的位置。}}{To the extent that this is true, it seems to me that legislatures are much better positioned to mount an assessment of the significance of an individual case in relation to a general issue of rights that affects millions and affects them in many different ways.}{承 K26，回應「法院面對具體個案」的優點（K67）：立法程序經遊說、聽證與辯論也考量個案，而且較能衡量個案與一般問\mbox{題的關係。}}
\keyentry{3}{K111}{1381}{}{從結果面的觀點看，讓權利分歧就權利法案的用語而爭出高下，是好主意還\mbox{是壞主意？}}{From an outcome-related point of view, is it a good idea or a bad idea that rights-disagreements be fought out in relation to the terms of a Bill of Rights?}{提出第三項優點的問題：權利法案的文本對爭點有何作用。K2\mbox{7 作答。}}
\keyentry{2}{K27}{1381}{}{權利法案所用的文字形式，當初想必不是顧及權利分歧而選定的。即使當初顧及了權利分歧，選定這些措辭想必也是為了迴避權利法案設立時既已存在的\mbox{權利分歧。}}{The forms of words used in the Bill of Rights will not have been chosen with rights-disagreements in mind. Or, if they were, they will have been chosen in order to finesse the disagreements about rights that existed at the time the Bill of Rights was set up.}{回應法院的第二項優點（依據文本）：權利法案的措辭並非為了處理權利分歧而選定，甚至是為了迴避制定當\mbox{時的分歧。}}
\keyentry{4}{K217}{1381}{四級・生動表述}{要使清醒、負責而出於善意的權利分歧探究有所聚焦，這些陳腔濫調式的表述可能恰恰\mbox{最不合適。}}{Their platitudes may be exactly the wrong formulations to focus clear-headed, responsible, and good faith explorations of rights-disagreements.}{直稱權利法案條文為「陳腔濫調」；接 K27\mbox{ 的論點。}}
\keyentry{3}{K112}{1381}{}{權利法案的書面表述也往往助長某種僵固的文本\mbox{形式主義。}}{The written formulations of a Bill of Rights also tend to encourage a certain rigid textual formalism.}{文本的第二項缺點：條文措辭自成套語，法院依其行事方式必定執著於文本，議員則可不受條文措辭的限制。K2\mbox{8 承此。}}
\keyentry{4}{K218}{1381}{四級・生動表述}{美國憲政經驗的一個教訓是，每一項規定的文字往往自成生命，成為一句套語，人們關於該項權利想說的一切，都執意用\mbox{它來表達。}}{One lesson of American constitutional experience is that the words of each provision tend to take on a life of their own, becoming the obsessive catchphrase for expressing everything one might want to say about the right in question.}{條文文字「自成生命」、成為執念般的套語；導讀時可舉 due process、equal protectio\mbox{n 為例。}}
\keyentry{4}{K219}{1381}{四級・生動表述}{但為其推理尋求文本上的庇護所，是法院行事方式（modus operandi）的一部分，而法院必定傾向於相當執著地以權利法案的文\mbox{本為取向。}}{But it is part of the modus operandi of courts to seek textual havens for their reasoning, and they will certainly tend to orient themselves to the text of the Bill of Rights in a rather obsessive way.}{「文本上的庇護所」比喻法院躲進條文以避開直接的道德論證；接\mbox{ K28。}}
\keyentry{3}{K160}{1381}{}{美國的經驗證實了這一點：在多數司法意見書裡，關於解釋理論的論辯與對道德問題的直接論辯之間的比例是偏斜的，凡認為那些問題本身重要的人，都不可能認為這樣的比例\mbox{令人滿意。}}{American experience bears this out: The proportion of argument about theories of interpretation to direct argument about the moral issues is skewed in most judicial opinions in a way that no one who thinks the issues themselves are important can possibly regard as satisfactory.}{權利法案文本使法院的論辯偏向解釋理論，排擠道德問題的\mbox{直接論辯。}}
\keyentry{2}{K28}{1381}{}{由於法官（一如我們其餘的人）關切的是：准許他們決定這些問題的這套程序是否正當，他們便緊抓授權自己的文本、爭論其解釋，而不敢直接討論\mbox{道德理由。}}{Because judges (like the rest of us) are concerned about the legitimacy of a process that permits them to decide these issues, they cling to their authorizing texts and debate their interpretation rather than venturing out to discuss moral reasons directly.}{法官因在意自身決定的正當性而緊抓文本，不直接討論\mbox{道德理由。}}
\keyentry{3}{K68}{1382}{}{權利法案的文本可能扭曲司法推理，不只由於它納入了什麼，也由於它遺\mbox{漏了什麼。}}{The text of a Bill of Rights may distort judicial reasoning not only by what it includes but also by what it omits.}{權利法案未列積極權利，法官可能加重財產權與契約自由的分量，進而宣告補足這項欠缺的制\mbox{定法違憲。}}
\keyentry{2}{K29}{1382}{}{議員與法官一樣，會說明自己如此表決的理由。理由是在我們所稱的辯論中提出的，並刊載於英國國會議事錄（Hansard）或美國國會紀錄（Congressional Re\mbox{cord）。}}{Legislators give reasons for their votes just as judges do. The reasons are given in what we call debate and they are published in Hansard or the Congressional Record.}{回應法院的第三項優點（說明理由）：議員在辯論中說明理由，並刊載於\mbox{議事紀錄。}}
\keyentry{4}{K220}{1382}{四級・生動表述}{差別在於，律師受過訓練，會細讀法官所說明的理由；他們未受過細讀立法上之推理的訓練（儘管他們偶爾會為解釋的目的\mbox{而翻檢它）。}}{The difference is that lawyers are trained to close study of the reasons that judges give; they are not trained to close study of legislative reasoning (though they will occasionally ransack it for interpretive purposes).}{括號裡律師「翻箱倒櫃」般搜刮立法紀錄（ransack），點出立法說理被忽視的原因在讀\mbox{者的訓練。}}
\keyentry{2}{K113}{1383}{}{但若手上是一部兩三百年前的陳舊憲法，那麼所謂的說明理由就很可能是人為的\mbox{、扭曲的。}}{But if one has an antiquated constitution, two or three hundred years old, then the alleged reason-giving is likely to be artificial and distorted.}{回應「法院說明理由」的優點：K29 說議員同樣說明理由，本段轉談理由的品質，法院的理由很少是充分審議會檢視的理由，美國的說理多在連結十八、十九世\mbox{紀的文句。}}
\keyentry{4}{K221}{1383}{四級・生動表述}{在美國，所謂的「說明理由」通常是設法把法院所面對的決定，與某段十八或十九世紀未經深思的陳舊文句\mbox{連結起來。}}{In the United States, what is called ``reason-giving'' is usually an attempt to connect the decision the court is facing with some antique piece of ill-thought-through eighteenth- or nineteenth-century prose.}{把憲法條文稱為「未經深思的陳舊文句」，「說明理由」加引號；接 \mbox{K113。}}
\keyentry{4}{K222}{1383}{四級・生動表述}{（例如，爭論「實質正當程序」是不是自相矛盾的用語，是思考勞動法乃至墮胎權的最\mbox{佳框架嗎？）}}{(For example, is an argument about whether ``substantive due process'' is an oxymoron the best framework for thinking about labor law or, for that matter, abortion rights?)}{以反問嘲諷「實質正當程序」是否自相矛盾的爭論，竟成了思考勞動法與墮胎權的框架；可帶到 Lochner 與\mbox{ Roe。}}
\keyentry{4}{K223}{1383}{四級・生動表述}{法院說明理由時，還要設法在眼前的決定與其他決定之間硬找出類比或區別；那些決定只是恰好先前提交到它們面前，而它們在那些決定裡也有過類似的\mbox{牽強推演。}}{Courts' reason-giving also involves attempts to construct desperate analogies or disanalogies between the present decision they face and other decisions that happen to have come before them (and in which they were engaged in similar contortions).}{「走投無路的類比」與「扭曲身段」（contortions），把援引先例寫成勉強的\mbox{肢體動作。}}
\keyentry{4}{K224}{1383}{四級・生動表述}{而在這整個過程裡，出於善意的權利分歧中所涉的真正問題被推到了邊緣。它們通常只占一份二十頁或更長的意見書裡的一兩段，而即使在那一兩段裡，這些問題也很少被\mbox{直接處理。}}{And all the time, the real issues at stake in the good faith disagreement about rights get pushed to the margins. They usually take up only a paragraph or two of the twenty pages or more devoted to an opinion, and even then the issues are seldom addressed directly.}{「二十頁裡的一兩段」的篇幅對比，緊接 K161 的 Roe 案五十頁；三\mbox{句可連讀。}}
\keyentry{4}{K161}{1383}{四級・生動表述}{例如，在聯邦最高法院 Roe v. Wade 一案五十頁的意見書裡，只有兩三段處理生育權在隱私上的道德重要性，而所涉的另一個道德問題是胎兒的權利地位，處理此一問題的那幾段，大部分用於呈現各方意\mbox{見的紛紜。}}{In the Supreme Court's fifty-page opinion in Roe v. Wade, for example, there are but a couple of paragraphs dealing with the moral importance of reproductive rights in relation to privacy, and the few paragraphs addressed to the other moral issue at stake—the rights-status of the fetus—are mostly taken up with showing the diversity of opinions on the issue.}{Roe 案與 1966 年下議院辯論對比的前半：意見書處理道德問\mbox{題的篇幅。}}
\keyentry{4}{K30}{1383}{四級・生動表述}{讀一讀那幾段：結論或許吸引人，但那「推理」\mbox{是單薄的。}}{Read those paragraphs: The result may be appealing, but the ``reasoning'' is thread-bare.}{作者對 Roe v. Wade 說理的評語，與 K31 的英國\mbox{辯論對照。}}
\keyentry{3}{K114}{1384}{}{但在關於違憲審查是否優於立法的結果面論證裡，這一點對法院\mbox{極為不利。}}{But it counts heavily against the courts in the outcome-related argument about the preferability of judicial review over legislation.}{法院關切自身的正當性，說理因而集中於權限、文本與先例；議員至少往往直接切入事情的核心。與 K3\mbox{0 並讀。}}
\keyentry{4}{K225}{1384}{四級・生動表述}{法院被這些正當性問題分散了注意力，於是聚焦於其他法院做過什麼、或權利法案的文字是什麼，而議員儘管有種種缺點，至少往往直接切入事\mbox{情的核心。}}{Distracted by these issues of legitimacy, courts focus on what other courts have done, or what the language of the Bill of Rights is, whereas legislators—for all their vices—tend at least to go directly to the heart of the matter.}{法院與議員的對照，插入語「儘管有種種缺點」先讓一\mbox{步再翻轉。}}
\keyentry{3}{K115}{1384}{}{就此而言，在沒有違憲審查的國家，立法機關辯論重要權利問題時，推理之豐富\mbox{令人矚目。}}{In this regard, it is striking how rich the reasoning is in legislative debates on important issues of rights in countries without judicial review.}{引入 1966 年英國下議院墮胎法案二讀辯論的例子；反對一方也稱道其立場受到尊重，與 Roe v. Wade 對照。K3\mbox{1 承此。}}
\keyentry{4}{K226}{1384}{四級・生動表述}{要找政治機關正面處理道德問題的例證，該法案的二讀辯論是最好不\mbox{過的一個。}}{The second reading debate on that bill is as fine an example of a political institution grappling with moral issues as you could hope to find.}{作者推舉 1966 年墮胎法案二讀為立法審議的範例，是全文對英國下議院最高的評價；接 K191、\mbox{K192。}}
\keyentry{4}{K227}{1384}{四級・生動表述}{那是一場持續的辯論，在英國國會議事錄裡約\mbox{占一百頁}}{It is a sustained debate—about one hundred pages}{「議事錄一百頁」與 K161 Roe 案意見書五十頁裡只有兩三段談道德問題，構成全文最常被記住的\mbox{篇幅對比。}}
\keyentry{4}{K191}{1384}{四級・生動表述}{他們熱切地辯論這些問題，但也辯得徹底而光明磊落，兼顧兩方的權利、原則與\mbox{務實問題。}}{They debated the questions passionately, but also thoroughly and honorably, with attention to the rights, principles, and pragmatic issues on both sides.}{1966 年下議院墮胎法案二讀辯論的描述：約一百頁、跨黨派、直接處理道德問題，與 Roe 案意見書（K16\mbox{1）對照。}}
\keyentry{2}{K192}{1384}{}{有一件事很不尋常：參與辯論的每一個人，甚至包括反對墮胎立場的下議院議員（當他們看出表決將往哪個方向走時），都稱道自己的立場在該討論中受到尊重、被\mbox{認真聆聽。}}{One remarkable thing was that everyone who participated in the debate, even the pro-life MPs (when they saw which way the vote was going to go), paid tribute to the respectfulness with which their positions had been listened to and heard in that discussion.}{對比的關鍵事實：落敗一方仍稱許受到尊重；下一句 K162 反問 Roe 案中反墮胎一方\mbox{何曾如此。}}
\keyentry{4}{K162}{1385}{四級・生動表述}{想一想：以聯邦最高法院審理 Roe v. Wade 為例，我們何曾聽過反對墮胎的一方稱道自己的立場在該案中受到認真而尊\mbox{重的討論？}}{Think about that: How many times have we ever heard anybody on the pro-life side pay tribute to the attention and respectfulness with which her position was discussed, say, by the Supreme Court in Roe v. Wade?}{對比的後半：下議院辯論中連反對陣營也稱許受到尊重，Roe\mbox{ 案沒有。}}
\keyentry{4}{K228}{1385}{四級・生動表述}{在美國，我們把墮胎權這類個人權利問題交付法院作憲法審判，並為此沾沾自喜，理由是法院可以視為「原則的論壇」，這是 Ronald Dworkin\mbox{ 的名言。}}{In the United States, we congratulate ourselves on consigning issues of individual rights such as abortion rights to the courts for constitutional adjudication on the ground that courts may be regarded as forums of principle, to use Ronald Dworkin's famous phrase.}{「沾沾自喜」的自嘲，加上 Dworkin「原則的論壇」一語；導讀時接 K69 對這句\mbox{話的反駁。}}
\keyentry{4}{K229}{1385}{四級・生動表述}{我們甚至有時說英國人不這樣做\mbox{是落後的。}}{Indeed we sometimes say the British are backward for not doing things that way.}{美國人笑英國落後，作者隨即以下議院辯論反證；自嘲美國\mbox{的優越感。}}
\keyentry{3}{K31}{1385}{}{但英國的立法辯論與美國的司法推理之間的關鍵差別在於，後者大多關切解釋與法理，而在前者，作決定的人能夠堅定地聚焦於墮胎這個問題本身及其所牽\mbox{涉的事項}}{But the key difference between the British legislative debate and the American judicial reasoning is that the latter is mostly concerned with interpretation and doctrine, while in the former decisionmakers are able to focus steadfastly on the issue of abortion itself and what it entails}{英國 1966 年墮胎法案的國會辯論與美國司法說\mbox{理的比較。}}
\keyentry{3}{K163}{1386}{}{確實，正如法院處理權利問題的方式會扭曲真正所涉的事項，當立法機關的多數出於恐慌行事、輕率行事，或在其偽辯論中一味鸚鵡學舌，複誦流行的或派系的口號時，立法上的推理同樣可能\mbox{令人蒙羞。}}{It is certainly the case that just as courts address questions of rights in ways that distort what is really at stake, so too can legislative reasoning be a disgrace, as legislative majorities act out of panic, recklessly, or simply parrot popular or sectarian slogans in their pseudo-debates.}{作者承認立法的推理也有缺陷，下一句接著問兩個機關的缺陷何者屬常態（K69\mbox{、K70）。}}
\keyentry{2}{K69}{1386}{}{儘管 Dworkin 把法院標榜為「原則的論壇」，我認為法院本就被期待以我所批評的那些方式行事：聚焦於先例、文本、法理與其他法條主\mbox{義的東西。}}{Despite Dworkin's rhetoric about ``forums of principle,'' I think courts are expected to behave in the ways that I have criticized, focusing on precedent, text, doctrine, and other legalisms.}{同段先問兩種機關審議上的缺陷何者屬常態、何者屬偏差；此句回答法院一方：聚焦先例、文本與法理，正是法\mbox{院的常態。}}
\keyentry{2}{K70}{1386}{}{然而就立法機關而言，草率或出於派系的立法，並不在我們通常所理解的立法機關設置\mbox{目的之內。}}{In the case of legislatures, however, hasty or sectarian legislating is not part of the normal theory of what legislatures are set up to do.}{回答立法機關一方：草率或派系的立法不屬立法機關的常態，核心情形不作此假設；與 K6\mbox{9 並讀。}}
\keyentry{4}{K230}{1386}{四級・生動表述}{或許有些國家（可能包括美國）形成了該國特有的立法弊病。若是如此，美國人就應當把他們支持違憲審查的非核心論證限於他們自身的\mbox{例外情況。}}{There may be some countries—perhaps the United States—in which peculiar legislative pathologies have developed. If that is so, then Americans should confine their non-core argument for judicial review to their own exceptional circumstances.}{插入語「可能包括美國」的揶揄；把美國的病例歸為非核心情形，預告\mbox{第柒部分。}}
\subsection*{伍、程序面的理由}
\keyentry{4}{K231}{1386}{四級・生動表述}{這個制度的辯護者編造了一些薄弱的程序面論證，我將在本部分末尾\mbox{檢視它們。}}{Some feeble process-related arguments have been concocted by defenders of the practice, and I shall review those at the end of this Part.}{「編造」「薄弱」兩詞先替本部分末尾的五個辯護下了評價；導讀時可預\mbox{告那五點。}}
\keyentry{3}{K71}{1386}{}{但大體如此：程序面的理由大多有利於\mbox{立法機關。}}{But it is mostly true: The preponderance of the process-related reasons weigh in favor of legislatures.}{第伍部分的結論先行；辯護者所提的程序面論證，在本部分末尾的五項調和論\mbox{證裡檢視。}}
\keyentry{4}{K232}{1386}{四級・生動表述}{（恐怕以下會\mbox{相當抽象。）}}{(I am afraid this is going to be quite abstract.)}{作者自承接下來的正當性理論很抽象，括號裡向讀者打招呼；與 K193「翻白眼」同\mbox{一種口吻。}}
\keyentry{3}{K164}{1387}{}{但他們可能無法說服她，原因並非她遲鈍，只是 \textit{C}\textsubscript{\textit{n}} 對這個棘手而嚴重的事項（並非不合理地）始終\mbox{另有見解。}}{But they may fail, not because of any obtuseness on her part, but simply because \textit{C}\textsubscript{\textit{n}} continues (not unreasonably) to hold a different view on this vexed and serious matter.}{\textit{C}\textsubscript{\textit{n}} 的不同意並非不講理，呼應第四項假設的善意分歧；正當性因此須以\mbox{程序答覆。}}
\keyentry{2}{K116}{1387}{}{這種以程序為基礎的答覆，其理論就是政治正\mbox{當性理論。}}{The theory of such a process-based response is the theory of political legitimacy.}{程序面理由的定義：無法說服 \textit{C}\textsubscript{\textit{n}} 結果正確時，向她說明決定是以公平方式作成。K32 的兩個問題\mbox{由此提出。}}
\keyentry{3}{K32}{1387}{}{用 \textit{C}\textsubscript{\textit{n}} 可能的說法：「為什麼是他們？為什\mbox{麼不是我？」}}{As \textit{C}\textsubscript{\textit{n}} might put it, ``Why them? Why not me?''}{程序面理由的出發點：不同意決定的公民 \textit{C}\textsubscript{\textit{n}} 向決定者提出\mbox{的第一問。}}
\keyentry{3}{K165}{1387}{}{這個問題是一般性的，不只是 \textit{C}\textsubscript{\textit{n}} 個人乖僻的問題；因此 Co、Cp，以及享有特權的決定者以外的其他所有 C，若提出類似的問題，正當性理論也必須給出類\mbox{似的答覆。}}{Because the problem is general—it is not just a matter of \textit{C}\textsubscript{\textit{n}}'s idiosyncratic perversity—it will have to give a similar answer to similar questions from Co and Cp and all the other C's not included in the set of privileged decisionmakers.}{答覆必須一般化，對每一位被排除的公\mbox{民都適用。}}
\keyentry{4}{K233}{1387}{四級・生動表述}{現在讓我們把這套抽象的代數講\mbox{得具體些。}}{Let us now make this abstract algebra more concrete.}{自嘲 \textit{C}\textsubscript{\textit{n}}、C1…Cm 的符號是「代數」，接著換成五百名議員與兩億五千萬人（\mbox{K117）。}}
\keyentry{3}{K117}{1387}{}{她問：(1) 為什麼這一群大約五百位男女（立法機關的成員）應享有特權，去決定一個影響我與另外兩億五千萬人的權利問題？以及 (2) 即使我接受這五百人享有特權，有些議員與我看法相同，為什麼沒有給他們的見解更\mbox{大的分量？}}{She asks: (1) why should this bunch of roughly five hundred men and women (the members of the legislature) be privileged to decide a question of rights affecting me and a quarter billion others?; and (2) even if I accept the privileging of this five hundred, why wasn't greater weight given to the views of those legislators who agreed with me?}{K32 兩個問題對立法機關的具體版本，與 K72 對法院的\mbox{版本並讀。}}
\keyentry{2}{K166}{1387}{}{第一個問題的答覆，來自立法機關公平選舉的理論：決定由誰躋身參與此類決定的少數人之列、享有此一特權時，像 \textit{C}\textsubscript{\textit{n}} 這樣的人與其所有同胞受到\mbox{平等對待。}}{The answer to the first question is provided by the theory of fair elections to the legislature, elections in which people like \textit{C}\textsubscript{\textit{n}} were treated equally along with all their fellow citizens in determining who should be privileged to be among the small number participating in decisions of this kind.}{立法機關對 \textit{C}\textsubscript{\textit{n}} 第一問的答覆，與 K118（第二問的答覆）合為立法機關的\mbox{兩項答覆。}}
\keyentry{2}{K118}{1388}{}{MD 比任何其他規則都更能在相互爭執的結果之間保持中立、平等對待參與者，並在所有意見分量均等的前提下，給每一個表達出來的意見盡可能\mbox{大的分量。}}{Better than any other rule, MD is neutral as between the contested outcomes, treats participants equally, and gives each expressed opinion the greatest weight possible compatible with giving equal weight to all opinions.}{立法機關對第二個問題的答覆：多數決的公平性。對第一個問題的答覆是公平選舉，見\mbox{同段前半。}}
\keyentry{3}{K119}{1388}{}{其理論是：兩者合起來，合理地近似於在全體公民之間以 MD 為決定程序（因而也合理地近似於把支撐 MD 的價值適用於\mbox{全體公民）。}}{The theory is that together these provide a reasonable approximation of the use of MD as a decision-procedure among the citizenry as a whole (and so a reasonable approximation of the application of the values underlying MD to the citizenry as a whole).}{選舉代表與代表之間的表決都用多數決，兩個答覆前後一貫；法院的情形缺少這種一貫，見\mbox{ K37。}}
\keyentry{2}{K120}{1388--1389}{}{在其他每一個人都享有平等發言權的前提下，我們給每一個人盡可能大的發言權。這就是我\mbox{們的原則。}}{We give each person the greatest say possible compatible with an equal say for each of the others. That is our principle.}{對 \textit{C}\textsubscript{\textit{n}} 的答覆歸結為政治平等的原則；P44 據此說法院的程序提不出相\mbox{應的答覆。}}
\keyentry{3}{K121}{1389}{}{但一套好的正當性理論（就現實世界的政體而言）會保留某種寬鬆，以容納不可避\mbox{免的缺陷。}}{But a good theory of legitimacy (for real-world polities) will have a certain looseness to accommodate inevitable defects.}{選舉與立法程序的缺陷可以承認，第一項假設只要求合理的公平；失能已無從矯正者屬非核心情形。K3\mbox{3 承此。}}
\keyentry{3}{K167}{1389}{}{它談的是合理的公平，不是完\mbox{美的公平。}}{It will talk about reasonable fairness, not perfect fairness.}{現實制度的缺陷由第一項假設吸收，正當性理論只要求合理的公平；緊接 \mbox{K121。}}
\keyentry{2}{K33}{1389}{}{正當性有一部分要經由比較\mbox{才能判斷。}}{Legitimacy is partly comparative.}{正當性有一部分須經比較判斷，下文據此比較立法與司法的正\mbox{當性基礎。}}
\keyentry{3}{K72}{1390}{}{她問：(1) 為什麼應由這九位男女來決定該事項？以及 (2) 即使由他們決定，為什麼他們採用現行的程序作決定，而不採用另一種程序，讓持 \textit{C}\textsubscript{\textit{n}} 所贊同之見解的大法官有更\mbox{大的分量？}}{She asks: (1) why should these nine men and women determine the matter?; and (2) even if they do, why should they make their decision using the procedure that they use rather than a procedure that gives more weight to Justices with a view that \textit{C}\textsubscript{\textit{n}} favors?}{\textit{C}\textsubscript{\textit{n}} 向立法機關提出的兩個問題（K32）改向法院提出，\mbox{兩相對照。}}
\keyentry{2}{K73}{1390}{}{這些問題，該院答覆起來遠比\mbox{議員困難。}}{These are much tougher questions for the Court to answer than they were for legislators to answer.}{法院答覆這兩個問題遠比立法機關困難；作者隨後長篇引 Scalia 的不同\mbox{意見為證。}}
\keyentry{4}{K234}{1390}{四級・生動表述}{我們有可靠的根據相信，這類質疑常在該院之外喧騰，而大法官有時因此感到苦惱。然而其中有些人會反思\mbox{那種苦惱。}}{We have it on good authority that challenges like these are often voiced noisily outside the Court and that the Justices are sometimes distressed by them. Some of them, however, reflect on that distress.}{「有可靠的根據」是反語，根據就是下面引的 Scalia 不同意見；接 \mbox{K193。}}
\keyentry{4}{K193}{1390}{四級・生動表述}{（讀者現在可以翻白眼，接下來幾分鐘不必理會，因為我要引用 Antonin Scalia 大法官，而且要長\mbox{篇引用他。）}}{(It is time to roll your eyes now and pay no attention for a few minutes, because I am going to quote Justice Antonin Scalia and quote him at length.)}{作者自嘲要長引 Scalia；引文用以說明法院的正當性問題\mbox{首當其衝。}}
\keyentry{4}{K235}{1390}{四級・生動表述}{說實在的，對於⋯⋯針對本院而來的「政治壓力」，我與本院一樣苦惱：那些遊行、那些信件、那些意圖促使我們改變意\mbox{見的抗議。}}{In truth, I am as distressed as the Court is … about the ``political pressure'' directed to the Court: the marches, the mail, the protests aimed at inducing us to change our opinions.}{Scalia 列舉遊行、信件、抗議三樣東西，畫出法院門外的景象；作者借對手陣營大法官之口說出正\mbox{當性問題。}}
\keyentry{3}{K122}{1390}{}{那個原因瀰漫於今日的意見書之中：一種新的憲法審判方式：認定法律時，憑藉的不是文本與傳統的做法，而是本院所稱的「說理的判斷」，而那到頭來不過是哲學上的偏好與道德\mbox{上的直覺。}}{That cause permeates today's opinion: a new mode of constitutional adjudication that relies not upon text and traditional practice to determine the law, but upon what the Court calls ``reasoned judgment,'' which turns out to be nothing but philosophical predilection and moral intuition.}{本段全為 Scalia 在 Casey 案的不同意見：大法官作價值判斷，人民便以政治壓力相待。作者以此說明 K73 的兩個問題對法院\mbox{更難回答。}}
\keyentry{4}{K236}{1390}{四級・生動表述}{這一切與針對本院施加「政治壓力」這件惱人的事相關，原因是兩項相伴的事實：美國人民熱愛民主，而美國人民\mbox{不是傻子。}}{What makes all this relevant to the bothersome application of ``political pressure'' against the Court are the twin facts that the American people love democracy and the American people are not fools.}{Scalia 的名句「美國人民不是傻子」；導讀時可問：為何一位保守派大法官的話正好支持 Wal\mbox{dron。}}
\keyentry{4}{K237}{1390}{四級・生動表述}{文本與傳統是供研究的事實，不是供人為之示\mbox{威的信念。}}{Texts and traditions are facts to study, not convictions to demonstrate about.}{「研究」與「示威」的對比，Scalia 以此區分律師的工作與\mbox{價值判斷。}}
\keyentry{3}{K123}{1391}{}{人民知道，他們的價值判斷與任何法學院所教的一樣好——也\mbox{許還更好。}}{The people know that their value judgments are quite as good as those taught in any law school—maybe better.}{Scalia 引文的後半，與 P41 同\mbox{一段意見。}}
\keyentry{3}{K124}{1391}{}{一個答覆可能是：大法官由總統與參議院任命並同意，而這兩個決定者與決定機關都具有某種來自選舉\mbox{的正當性。}}{One answer might be that the Justices have been appointed and approved by decisionmakers and decisionmaking bodies (the President and the Senate) who have certain elective credentials.}{法院對第一個問題的可能答覆。K34 比較兩種選任方式，結論是立法\mbox{選舉較優。}}
\keyentry{4}{K238}{1391}{四級・生動表述}{所以，面對我們這位公民的質疑，違憲審查的辯護者並非完全語塞，仍有話可說。然而，正當性若須經由比較判斷，這個答覆就\mbox{極其不足。}}{So, the defender of judicial review is not altogether tongue-tied in response to our citizen's challenge; there is something to say. Nevertheless, if legitimacy is a comparative matter, then it is a staggeringly inadequate response.}{先讓步「並非完全語塞」，再下「極其不足」的重話；導讀時接\mbox{ K34。}}
\keyentry{2}{K34}{1391}{}{立法機關的選舉制度也不完美，但就民主與民主價值而言，它顯然優於司法機關那種間接而有限的民主正\mbox{當性基礎。}}{The system of legislative elections is not perfect either, but it is evidently superior as a matter of democracy and democratic values to the indirect and limited basis of democratic legitimacy for the judiciary.}{立法選舉與司法任命的民主正\mbox{當性比較。}}
\keyentry{4}{K239}{1391}{四級・生動表述}{法院以表決作成決定，把 MD 原則適用於寥寥數人，我對此一直\mbox{深感好奇。}}{I have always been intrigued by the fact that courts make their decisions by voting, applying the MD principle to their meager numbers.}{「深感好奇」是反語，「寥寥數人」對照立法機關五百人；接 K35 的\mbox{「數人頭」。}}
\keyentry{2}{K35}{1391}{}{但到頭來歸結為數人頭：在美國聯邦最高法院，五票擊敗四票，不論大法官編造出\mbox{什麼論證。}}{But in the end it comes down to head-counting: five votes defeat four in the U.S. Supreme Court, irrespective of the arguments that the Justices have concocted.}{回應「法院經由說理決定」：法院內部最後同樣以多\mbox{數票決定。}}
\keyentry{2}{K36}{1392}{}{他們不代表任何人。他們得以參與，所憑的是職能，而不是應\mbox{享的資格。}}{They do not represent anybody. Their claim to participate is functional, not a matter of entitlement.}{法官參與決定的依據在於職能，與議員代表\mbox{選民不同。}}
\keyentry{2}{K37}{1393}{}{立法機關使用 MD，另有公平性的論證支持；法院使用 MD，則沒有這\mbox{樣的論證。}}{There is no additional fairness argument for the use of MD by courts, as there is for its use by legislatures.}{多數決的公平性論據只見於立法機關，法院使用多數決時沒有這一論據。第伍部分對 K35 的\mbox{推論落點。}}
\keyentry{2}{K125}{1393}{}{但我認為我已能在本部分與前一部分顯示：結果面的論據不具決定性（或者它有利於立法機關），而程序面的論據幾乎全在立法機\mbox{關這一方。}}{But I think I have been able to show in this Part, and the previous Part, that the outcome-related case is inconclusive (or it argues in favor of legislatures) while the process-related case is almost wholly on the legislative side.}{第肆、伍兩部分的小結，與 K22、K23 的預告相應；同段說法院的程序對政治平等之原則提\mbox{不出答覆。}}
\keyentry{4}{K240}{1393}{四級・生動表述}{違憲審查的辯護者佯稱有這樣的論據。但如我們上面所見，那只是未經支\mbox{持的斷言。}}{Defenders of judicial review pretend that there is. But as we saw above, it is just unsupported assertion.}{「佯稱」與「未經支持的斷言」，第肆、伍部分結\mbox{論的重話。}}
\keyentry{4}{K241}{1393}{四級・生動表述}{或許是察覺到這一切，違憲審查的辯護者作了若干最後一搏的嘗試，想使他們所偏好的制度與民主價值相調和。我將簡短地檢視這些嘗試，因為它們沒有多少\mbox{實質內容。}}{Perhaps aware of all this, defenders of judicial review have tried a number of last-ditch attempts to reconcile their favored institution to democratic values. I will consider these briefly, because there is not much to them.}{「最後一搏」（last-ditch，壕溝戰的比喻）與「沒有多少實質內容」；導讀時用來帶出其後\mbox{五點辯護。}}
\keyentry{2}{K74}{1393}{}{權利法案的抽象用語，充其量是人民選定來爭執這些問題的場域。我們一直在考慮的問題是：在那些場域上爭出高下的問題，由\mbox{誰來解決。}}{At most, the abstract terms of the Bill of Rights are popularly selected sites for disputes about these issues. The question we have been considering is who is to settle the issues that are fought out on those sites.}{回應第一項調和論證（法官只是執行權利法案所體現的人民決定）：權利法案只提供爭執的場域，問題仍在由誰解決。與 K1\mbox{9 合讀。}}
\keyentry{4}{K194}{1393}{四級・生動表述}{第二，本於大致相同的精神，違憲審查的辯護者主張，法官只是在執行該社會自己對權利的預先承諾。依這種說法，該社會已就某些權利原則把自己綁在桅桿上，而法官就像 Ulysses 的船員，只是確保繩索\mbox{沒有鬆開。}}{Second, and in much the same spirit, defenders of judicial review claim that judges are simply enforcing the society's own precommitment to rights. The society has bound itself to the mast on certain principles of right, and, like Ulysses' shipmates, the judges are just making sure the ropes remain tied.}{Ulysses 之喻：第二項調和論證的內容；作者的反駁見 \mbox{K126。}}
\keyentry{2}{K126}{1394}{}{Ulysses 模型只在預先承諾用以防範各種反常時才能成立；人們對何為合理結果確有真正分歧時，預先承諾若用以阻止人們改變主意，這個模型就\mbox{不能成立。}}{The Ulysses model works only when the precommitment guards against various aberrations, not when it guards against changes of mind in relation to genuine disagreement as to what a reasonable outcome would be.}{對辯護者第二項調和論證（預先承諾）的回應：社會並未承諾某一權利的\mbox{特定見解。}}
\keyentry{2}{K127}{1394}{}{這個論證有缺陷，因為它並未為司法決定所確立的那個基準線\mbox{提出辯護。}}{This argument is flawed because it does not defend the baseline that judicial decisionmaking establishes.}{對第三項調和論證（未修憲即默示認可）的回應：修憲須特別多數，偏向司法決定的基準線本身\mbox{仍待證立。}}
\keyentry{3}{K147}{1394}{}{修改權利法案通常需要特別多數；或者，如果它是英國或紐西蘭那種形式的制定法，它在政治文化中會具有一定的地位，從而抬高修法的代價、加重修\mbox{法的負擔。}}{Amending a Bill of Rights characteristically involves a supermajority; or if it is a British- or New Zealand-style statute, it will have credentials in the political culture that raise the stakes and increase the burden associated with the amendment effort.}{K127 的理由之一：推翻判決須過特別多數或相當的政治門檻，現狀卻由\mbox{判決定下。}}
\keyentry{3}{K148}{1394}{}{如果我們那位不滿的公民 \textit{C}\textsubscript{\textit{n}} 問為什麼要這樣事先做手腳，讓局面偏向一方，我們唯一能給她的答覆又回到\mbox{司法決定。}}{If our disgruntled citizen \textit{C}\textsubscript{\textit{n}} asks why the deck should be stacked in this way, the only answer we can give her refers back to judicial decision.}{K127 的理由之二：門檻偏向一方的理由只能回到判決本身，接回 \textit{C}\textsubscript{\textit{n}}\mbox{ 的問題。}}
\keyentry{2}{K128}{1394}{}{此言不虛；但（如我已經指出的）這個問題須經比較才能判斷，而法官這種民主正當性的根據，根本無法與民選議\mbox{員的相比。}}{This is true; but (as I have already remarked) the issue is comparative, and these credentials are not remotely competitive with the democratic credentials of elected legislators.}{對第四項調和論證（法官的民主正當性）的回應，承 K34；同段指出此說也削弱違憲審查別具價值的\mbox{積極論據。}}
\keyentry{2}{K168}{1394}{}{此外，我們越是因為法官有民主正當性的根據而接受他們，就越是動搖支持違憲審查的積極論據；該論據主張違憲審查是一種別具價值的政治\mbox{決定形式。}}{Moreover, to the extent that we accept judges because of their democratic credentials, we undermine the affirmative case that is made in favor of judicial review as a distinctively valuable form of political decisionmaking.}{對第四項調和論證的第二層反駁：以民主資格為法官辯護，等於進入法官必輸的程\mbox{序面比較。}}
\keyentry{2}{K75}{1395}{}{但把違憲審查嵌入一組更廣的公民參與途徑之中，並不改變以下事實：政治平等原則通常被視為民主的根本，而這種公民參與途徑並不\mbox{受其約束。}}{But embedding judicial review in a wider array of modes of citizen participation does not alter the fact that this is a mode of citizen involvement that is undisciplined by the principles of political equality usually thought crucial to democracy.}{回應第五項調和論證（違憲審查是公民參與的另一條途徑）：這條途徑不受政治平等\mbox{原則約束。}}
\keyentry{4}{K242}{1395}{四級・生動表述}{人們若希望自己的意見得到比選舉政治所給予者更大的分量，往往求諸\mbox{違憲審查。}}{People tend to look to judicial review when they want greater weight for their opinions than electoral politics would give them.}{一句點破訴訟者的動機：想要比一人一票更大的分量；導讀時可接\mbox{政治平等。}}
\subsection*{陸、多數暴政}
\keyentry{4}{K243}{1395}{四級・生動表述}{就核心情形而言，我要給違憲審查的辯護者最後一次嘗\mbox{試的機會。}}{I want to give defenders of judicial review—for the core case—one last bite at the apple.}{「再咬一口蘋果」的口語（再給一次機會），譯文未保留字面；引出多\mbox{數暴政論。}}
\keyentry{3}{K129}{1395}{}{關於民主立法機關的工作，人們最常表達的憂慮是：由於立法機關以多數主義為基礎組成，立法程序可能成為「多數暴政\mbox{」的表現。}}{The concern most commonly expressed about the work of a democratic legislature is that, because they are organized on a majoritarian basis, legislative procedures may give expression to the ``tyranny of the majority.''}{第陸部分處理的反對論點。K38 說此\mbox{論點混淆。}}
\keyentry{4}{K244}{1395}{四級・生動表述}{這種恐懼如此普遍，在我們的政治文化中又如此司空見慣，「多數暴政」這個說法如此輕易就掛在\mbox{我們嘴邊}}{So widespread is this fear, so familiar an element is it in our political culture, so easily does the phrase ``tyranny of the majority'' roll off our tongues}{三個「so…」排比，寫「多數暴政」成了脫口而出的套語；導讀時用來說明作者要拆解的是這\mbox{個慣用語。}}
\keyentry{3}{K38}{1395}{}{我認為這個常見的論證有嚴\mbox{重的混淆。}}{I believe that this common argument is seriously confused.}{第陸部分開頭對多數暴政論\mbox{證的總評。}}
\keyentry{4}{K195}{1396}{四級・生動表述}{例如，吸用佩奧特（peyote）者會認為，把他們的聖禮置於普遍適用的麻醉藥品法之\mbox{下是暴政。}}{For example, the peyote smokers will think the subjection of their sacraments to generally applicable narcotics laws is tyrannical.}{暴政之說的例子：任何權利分歧的一方都可能指對\mbox{方為暴政。}}
\keyentry{3}{K169}{1396}{}{的確，在有些情形裡，一個權利問題的兩方都會提出暴\mbox{政的指控。}}{Indeed, in some cases, there will be allegations of tyranny on both sides of a rights issue.}{暴政之說兩方都可援用，因此不能直接支持\mbox{違憲審查。}}
\keyentry{2}{K130}{1396}{}{但任何決定程序都是如此。法院有時也會\mbox{施行暴政。}}{But the same is true of any decision process. Courts will sometimes act tyrannically as well.}{承 P21：錯誤的權利決定在任何程序都會發生，多寡的比較已在第肆\mbox{部分處理。}}
\keyentry{4}{K245}{1396}{四級・生動表述}{我先撇開一個吹毛求疵的論點不談：法院也可能以多數表決作\mbox{成其決定。}}{I leave aside the pedantic point that a court may also reach its decision by majority voting.}{說要撇開，其實順帶點出：法院同樣靠多數決；自嘲「\mbox{吹毛求疵」。}}
\keyentry{2}{K131}{1396}{}{由民眾多數所施行的暴政（例如由民選代表的多數，而其中每一位都得到其選區多數的支持），是一種特別惡劣的暴政形態嗎？我看不出它\mbox{怎麼會是。}}{Is tyranny by a popular majority (e.g., a majority of elected representatives, each supported by a majority of his constituents) a particularly egregious form of tyranny? I do not see how it could be.}{作者的主張：由多數作成並不加重暴政；同段並主張多數主義的面向反而減輕暴政，因為決定的作成未排除任何人以平等者的\mbox{身分參與。}}
\keyentry{2}{K170}{1396}{}{我們或者說，不論暴政性的決定是如何作成、在誰之間作成，暴政就是暴政；或者說，多數主義這個面向其實減輕了暴政（這是我的見解），因為它顯示該決定至少有一點不屬暴政：它的作成並未以暴政的手段排除某些人以平等者的\mbox{身分參與。}}{Either we say that tyranny is tyranny irrespective of how (and among whom) the tyrannical decision is made, or we say—and this is my view—that the majoritarian aspect actually mitigates the tyranny, because it indicates that there was at least one non-tyrannical thing about the decision: It was not made in a way that tyrannically excluded certain people from participation as equals.}{回答「由多數加諸是否使暴政加重」：多數決反而減輕暴政，因為沒有人被排除於平等\mbox{參與之外。}}
\keyentry{4}{K246}{1396}{四級・生動表述}{這說法或許顯得有點輕佻，所以讓我以較不挑釁的方式處理\mbox{這個問題。}}{That may seem a little flip, so let me address the question less provocatively.}{作者自承前一段（K170）說得輕佻，再換個方式講；導讀時可提醒前後兩段是同一論點的\mbox{兩種說法。}}
\keyentry{3}{K132}{1397}{}{但在這個特定的脈絡裡，它們未必是與政治決定程序\mbox{有關的詞。}}{But in this particular context they are not necessarily terms related to political decision-processes.}{預告 K39 決定上的多數與議題上的多\mbox{數的區分。}}
\keyentry{2}{K39}{1397}{}{然而，當它發生時，我們至少要先把「決定上的」（decisional）多數與少數，和我所稱「議題上的」（topical）多數與少數\mbox{區分開來}}{When it does, however, we need to distinguish at least in the first instance between the ``decisional'' majority and minority and what I shall call the ``topical'' majority and minority}{決定上與議題上的多數與少數之分，第陸部分的\mbox{分析工具。}}
\keyentry{4}{K247}{1397}{四級・生動表述}{例如種族不正義往往就是如此：白人議員（決定上的多數）投票支持白人特權（議題上的多數）；黑人議員在爭取黑人平等權利的奮\mbox{鬥中落敗。}}{This is often true in the case of racial injustice for example: White legislators (decisional majority) vote for white privilege (topical majority); black legislators lose out in the struggle for equal rights for blacks.}{用白人議員與黑人議員的畫面，具體說明決定上的多數與議題上的多數重合（K39\mbox{、K40）。}}
\keyentry{3}{K133}{1397}{}{但到此為止，還沒有證明我遭受了\mbox{任何暴政。}}{But so far it has not been shown that anything tyrannical has happened to me.}{以投票落敗的公民為例套用 K39 的區分：證明暴政還須證明決定錯誤，且其屬於受不利影響的議題\mbox{上的少數。}}
\keyentry{2}{K171}{1397}{}{要證明此事，還須證明另外兩件事：(1) 該決定就其對受影響者權利的意涵而言確實是錯誤的、屬於暴政的；以及 (2) 我屬於議題上的少數，而該少數的權利受到該錯誤決定的\mbox{不利影響。}}{To show that we would have to show two additional things: (1) that the decision really was wrong and tyrannical in its implications for the rights of those affected; and (2) that I was a member of the topical minority whose rights were adversely affected by this wrong decision.}{表決落敗構成暴政的兩項要件；決定上的少數與議題上的少數之分由此推出 K40\mbox{ 的界定。}}
\keyentry{2}{K134}{1398}{}{這裡要記住的是：我所屬的共同體沒有依我的意見行事，單憑此事，我並未遭受\mbox{任何暴政。}}{The point to remember here is that nothing tyrannical happens to me merely by virtue of the fact that my opinion is not acted upon by a community of which I am a member.}{屬於決定上的少數不等於受暴政；議題上的少數也可能對自己的權利判斷錯誤。K40 據此界定\mbox{多數暴政。}}
\keyentry{3}{K172}{1398}{}{人們（包括議題上的少數的成員）未必擁有他們自以為擁\mbox{有的權利。}}{People—including members of topical minorities—do not necessarily have the rights they think they have.}{即使屬於議題上的少數，也不當然\mbox{受到暴政。}}
\keyentry{2}{K40}{1398}{}{刻畫多數暴政最有成效的方式，是說它發生在議題上的少數與決定上的少數相\mbox{重合之時。}}{The most fruitful way of characterizing tyranny of the majority is to say that it happens when topical minorities are aligned with decisional minorities.}{以 K39 的區分界定多數暴政：議題上的少數同時是決定\mbox{上的少數。}}
\keyentry{2}{K135}{1398}{}{結論毋寧是：多數暴政——如果這個詞被負責任地使用——是非核心情形的一項特徵，在那種情形裡，人們對自己以外的少數者權利或個人權利\mbox{甚少關心。}}{Rather, the conclusion is that tyranny of the majority—if that term is being used responsibly—is a characteristic of non-core cases, in which people care little for minority or individual rights other than their own.}{多數暴政與第三項假設不相容，因而落在非核心情形，第柒\mbox{部分處理。}}
\keyentry{3}{K136}{1399}{}{他論稱，這不只是因為它可能產生暴政性的結果，更是因為對權利的尊重是任何政治決定體系之正當性的\mbox{背景條件。}}{This is not just because it may generate tyrannical outcomes, he argues, but because respect for rights is a background condition for the legitimacy of any system of political decisionmaking.}{Dworkin 的前提：對權利的尊重是多數決正當性的背景條件，每一位投票者須確信他人以平等的關切與尊重看待他。作者的回應見 K76\mbox{、K77。}}
\keyentry{4}{K196}{1399}{四級・生動表述}{一群恐怖分子以多數決決定我的命運，即使所用的 MD 程序也給了我一票，這樣的決定仍完全不具正當性，因為這項背景條件\mbox{並未滿足。}}{A bunch of terrorists deciding my fate by majority decision (even an MD process in which I am given a vote) has no legitimacy at all, because this background condition is not met.}{Dworkin 背景條件說的設例；作者的回應見 K76\mbox{、K77。}}
\keyentry{3}{K137}{1399}{}{Dworkin 說，答案完全取決於法院是否作成正\mbox{確的決定。}}{The answer, Dworkin says, depends entirely on whether the court makes the right decision.}{Dworkin 據其前提推論：法院宣告侵害民主條件的法律違憲，民主並無損失。作者的回應見\mbox{ K76。}}
\keyentry{2}{K76}{1400}{}{即使程序正是我們的分歧所涉者，我們仍然在意程序\mbox{上的價值。}}{We care about process-values even when process is what is at stake in our disagreements.}{回應 Dworkin 的第一點：決定的對象即使是民主本身，作成決定的程序仍須受程序\mbox{價值檢驗。}}
\keyentry{2}{K77}{1400}{}{但如果我們這樣讀它，那麼本文所考慮的那種社會便滿足 Dworkin\mbox{ 的前提。}}{But if we read it this way, then Dworkin's premise is satisfied for the sort of society we are considering in this Essay.}{作者稱為最有力的回應：Dworkin 的前提採第二種讀法（認真對待彼此的權利），本文所設想的社會即滿足之，多數決不因此失\mbox{其正當性。}}
\keyentry{3}{K138}{1400}{}{第二，我們也可以用本部分所發展的區分，來幫助處理以下這項指控：就權利所為的立法決定若不受審查，便是多數\mbox{自斷其案。}}{Second, we can also use the distinctions developed in this Part to help deal with the allegation that unreviewable legislative decisionmaking about rights involves the majority being the judge in its own case.}{第二項論證：援引 \textit{nemo iudex in sua causa} 的人主張關於權利的終局決定應移交法院。K41 與 P5\mbox{8 回應。}}
\keyentry{3}{K173}{1400}{}{幾乎任何想像得到的決定規則，最終都會涉及某人就其自身之\mbox{案作決定。}}{Almost any conceivable decision-rule will eventually involve someone deciding in his own case.}{回應 nemo iudex 的第一層：終局者總在裁斷自己的見解，法院也不例外；接\mbox{ K41。}}
\keyentry{4}{K248}{1400}{四級・生動表述}{除非我們設想一條字面意義上無盡的上訴鏈，否則總會有某個人或某個機關的決定\mbox{是終局的。}}{Unless we envisage a literally endless chain of appeals, there will always be some person or institution whose decision is final.}{「無盡的上訴鏈」的歸謬，說明任何決定程序最後都有人自己判自己的案；接 K173\mbox{、K41。}}
\keyentry{2}{K41}{1401}{}{輕率地援引 \textit{nemo iudex in sua causa}，不能成為忘記正當性之基本邏輯的藉口：人們有分歧，因而需要終局的決定與終局的\mbox{決定程序。}}{Facile invocations of \emph{nemo iudex in sua causa} are no excuse for forgetting the elementary logic of legitimacy: People disagree, and there is need for a final decision and a final decision-procedure.}{回應 \textit{nemo iudex in sua causa}：分歧存在，終局決定與終局決定程序\mbox{即不可少。}}
\keyentry{3}{K174}{1401}{}{但引人注意的是，這種情況極少發生，包括在美國通常由違憲審查處理的那幾類案件中也\mbox{極少發生。}}{But it is striking how rarely this happens, including how rarely it happens in the kinds of cases that are normally dealt with by judicial review in the United States.}{議題上的多數與決定上的多數重合甚少發生，下文舉墮胎與優惠性差別待遇為例；接 \mbox{K139。}}
\keyentry{2}{K139}{1401}{}{理應預期的是，認真對待權利的人會對權利有分歧；但這些分歧相對獨立於個人在該事項上的切身利害，正是他們認真對待權利的\mbox{一個表徵。}}{People who take rights seriously must be expected to disagree about them; but it is a sign of their taking rights seriously that these disagreements will be relatively independent of the personal stakes that individuals have in the matter.}{對「多數自斷其案」的回應：議題上的多數與決定上的多數很少重合，墮胎與優惠性差別待遇的支持者分布即是例證；重合普遍存在者屬非\mbox{核心情形。}}
\subsection*{柒、非核心情形}
\keyentry{3}{K78}{1401}{}{我認為，對許多人來說，支持違憲審查的論據建立在拒絕接受這些\mbox{假設之上。}}{For many people, I think the case for judicial review rests on the refusal to accept these assumptions.}{支持違憲審查者多以第一、第三項假設在現實中不成立為據，第柒部分\mbox{由此展開。}}
\keyentry{4}{K249}{1401}{四級・生動表述}{違憲審查有一部分是對人們所認知到的民主制度失靈的回應，或者有一部分是對以下事實的回應：許多人並未充分認真地對待權利（因此他們需要一個法院代\mbox{他們為之）。}}{Judicial review is in part a response to perceived failures of democratic institutions, or it is in part a response to the fact that many people do not take rights sufficiently seriously (so they need a court to do it for them).}{括號裡「需要一個法院代他們為之」，把違憲審查寫成替人民認真對待權利\mbox{的代理人。}}
\keyentry{3}{K175}{1402}{}{例如第三項假設（該社會一般信奉權利）相當容易滿足，因為支持違憲審查的論據幾乎總是假定：擬實行違憲審查的社會已有一部權利法案，而該法案與公民的見解有某種真\mbox{實的關聯。}}{Assumption three, for example—a general commitment to rights in the society—is fairly easily satisfied, given that the case for judicial review almost always assumes that somehow the society for which judicial review is envisaged has a Bill of Rights that stands in some real relation to the views of citizens.}{回應「假設不切實際」：第三項假設是違憲審查支持者自\mbox{己的前提。}}
\keyentry{2}{K79}{1402}{}{支持立法決定具有正當性的論據，並不依賴於任何關於立法制度臻於烏托邦式完美的假定，也不依賴於它們在選舉與程序面向上完美地體現政治平等原則。它取決於這些制度明白地以該原則為取向、以旨在滿足該原則的方式組成，並為此作出合\mbox{理的努力。}}{The case for the legitimacy of legislative decisionmaking does not depend on any assumption of the utopian perfection of legislative institutions, nor on their perfectly embodying the principle of political equality in their elective and procedural aspects. It turns on these institutions being explicitly oriented to this principle, organized in a way that is designed to satisfy the principle, and making a reasonable effort to do so.}{回應「假設不切實際」的批評：立法正當性的論據只要求制度以政治平等為取向、為此作合理努力，不\mbox{要求完美。}}
\keyentry{4}{K250}{1402}{四級・生動表述}{最後，我特意引用一個實際的立法機關的實際審議（英國下議院就 Medical Termination of Pregnancy Bill 1966 所為的審議）作為立法機關可能如何運作的例子，而不是憑哲學上的先驗（a priori）推想虛構出\mbox{來的例子。}}{Finally, I took care to cite the actual deliberations of an actual legislature—the House of Commons on the Medical Termination of Pregnancy Bill 1966—as an example of how legislatures might work, not some concoction of the philosophical a priori.}{「實際的立法機關的實際審議」對「哲學先驗的虛構」，回頭呼應第肆部分的下\mbox{議院辯論。}}
\keyentry{1}{K42}{1402}{一級・範圍}{在那些假設不成立的情形裡，本文所提出的反對違憲審查的論證\mbox{無法成立。}}{In cases in which the assumptions fail, the argument against judicial review presented in this Essay does not go through.}{回到 K09 的條件：假設不成立時，本文的論證不及\mbox{於該社會。}}
\keyentry{1}{K43}{1402}{}{然而，由此並不能推出：只要那些假設不成立，對法律的違憲審查就是可\mbox{資辯護的。}}{However, it does not follow that judicial review of legislation is defensible whenever the assumptions fail.}{假設不成立，並不當然支持違憲審查。讀者常把非核心情形讀成支持違憲審查的理由，此句即排除\mbox{這種讀法。}}
\keyentry{2}{K176}{1402}{}{如果一個社會的司法決定與其立法決定同樣腐化、同樣帶有偏見，那麼設立對法律的違憲審查可能\mbox{並不恰當。}}{It may not be appropriate to set up judicial review of legislation if judicial decisionmaking in a society is no less corrupt or no less prejudiced than its legislative decisionmaking.}{K43 的理由之一：比較的邏輯在非核心情形照樣適用，司法若同樣失靈，違憲審查不\mbox{當然成立。}}
\keyentry{3}{K80}{1403}{}{於是產生兩個問題：(1) 就立法機關而言，有沒有可能改善這個情況？(2) 假定法院會把那些問題處理得更好，重要權利問題的終局決定權是否應\mbox{交由法院？}}{Two questions then arise: (1) Is it possible to improve the situation, so far as the legislature is concerned? (2) Should a final power of decision for important issues of rights be vested in the courts, assuming that the courts would handle those issues better?}{第一項假設不成立時的兩個問題：立法機關能否改善，終局決定權應否\mbox{交給法院。}}
\keyentry{2}{K81}{1403}{}{把終局決定權交給法院，很可能使改革立法機關更為困難，或使第一項假設、或許還有第三項假設所預設的那種立法風氣\mbox{更難養成。}}{Vesting the final power of decision in courts may well make it more difficult to reform the legislature or more difficult to develop the legislative ethos that the first assumption, and perhaps also the third assumption, presuppose.}{K80 的兩個問題並非全然獨立：把終局決定權交給法院，可能阻礙立法機關的改革與立法風\mbox{氣的養成。}}
\keyentry{4}{K251}{1403}{四級・生動表述}{我聽過關於美國的這類推測：其想法是，美國的立法機關、特別是州立法機關，運作不負責任、未能認真對待權利，原因是議員知道有法院作後盾，負責任的文化因此難以在議員\mbox{之間養成。}}{I have heard speculation to this effect about the United States: The idea is that U.S. legislatures, particularly state legislatures, operate irresponsibly and in a way that fails to take rights seriously because the knowledge that the courts are there as backup makes it harder to develop a responsible culture among legislators.}{法院當「後盾」反而讓議員不負責任的推測；導讀時可討論違憲審查對立法文化的反作用（接\mbox{ K81）。}}
\keyentry{2}{K140}{1403}{}{在我看來，有一種非核心情形，支持對立法決定為違憲審查的論證在其中具有幾分說服力；用這段話來刻畫這種情形，\mbox{再好不過。}}{This it seems to me is an excellent way of characterizing the sort of non-core case in which the argument for judicial review of legislative decisions has some plausibility.}{作者承認 Carolene Products 案腳註四所述的偏見屬於第一項假設不成立的非核心情形，違憲審查在其中有幾分說服力。K43 限\mbox{定其範圍。}}
\keyentry{3}{K141}{1403}{}{從我在第陸部分所說的可以推出：並非每一個少數都值得這種特別待遇；當然不是每一個決定上的少數，甚至也不是每一個議題\mbox{上的少數。}}{It follows from what I said in Part VI that not every minority deserves this special treatment: certainly not every decisional minority, and not even every topical minority.}{以第陸部分的區分限定腳註四論證的\mbox{適用範圍。}}
\keyentry{3}{K177}{1403}{}{甚至沒有理由認為每一個長期的少數都值得這種特別待遇，更遑論長期的決定上的少數，例如美國的布\mbox{爾什維克。}}{There is no reason to suppose even that every chronic minority deserves this special treatment, certainly not chronic decisional minorities—Bolsheviks in the United States, for example.}{對 Carolene Products 腳註四的限縮：長期落敗者不當然適格；接 \mbox{K141。}}
\keyentry{3}{K142}{1404}{}{並非每一個界線分明、可資辨識的少數都是孤立\mbox{而隔絕的。}}{Not every distinct and identifiable minority is discrete and insular.}{「孤立而隔絕」指無法與他人結盟推進利益的少數，第陸部分所說兩種少數相重合\mbox{即屬此類。}}
\keyentry{4}{K252}{1404}{四級・生動表述}{Stone 大法官的用語本身沒有\mbox{什麼魔力。}}{There is nothing magical about Justice Stone's language.}{「沒有魔力」一語，去除 Carolene Products 腳註四「孤立而隔絕」的\mbox{神聖光環。}}
\keyentry{2}{K178}{1404}{}{但如果認真對待，「孤立」與「隔絕」是有用的形容詞，因為這兩個詞所傳達的，不只是一個獨立於政治決定而存在的少數（換句話說，議題上的少數），也是一個其成員與共同體其餘部分相隔絕的少數：所謂相隔絕，意思是他們與非成員之間沒有許多共同利益，足以使他們建立一系列聯盟來推進其利益。決定上的少數與議題上的少數相重合，即我們在第陸部分所警告的情形，是這個意義下「隔絕性」\mbox{的好例子。}}{But if taken seriously, ``discrete'' and ``insular'' are useful adjectives, for they convey not just the idea of a minority that exists apart from political decisionmaking—in other words a topical minority—but also a minority whose members are isolated from the rest of the community in the sense that they do not share many interests with nonmembers that would enable them to build a series of coalitions to promote their interests. The alignment of decisional and topical minorities that we warned against in Part VI is a good example of ``insularity'' in this sense.}{discrete 與 insular 的界定，並與第陸部分的兩種\mbox{少數相接。}}
\keyentry{3}{K143}{1404}{}{重點不在於堅持某一種特定形態的反感，而在於把各種形態的反感，與關於權利的合理分歧這個現象\mbox{區分開來。}}{The point is not to insist on any particular mode of antipathy, but to distinguish between its various modes and the phenomenon of reasonable disagreement about rights.}{腳註四的第二項判準（偏見）與第三、第四項假設不相容；作者要區分的是反感與合理分歧。K4\mbox{4 承此。}}
\keyentry{4}{K197}{1404}{四級・生動表述}{如果是，那麼這個情形可能不只是非核心的，而且\mbox{是無望的。}}{If they are, then the case may be not only non-core but hopeless.}{司法上的多數若與立法上的多數同樣帶有偏見，違憲審查亦\mbox{無濟於事。}}
\keyentry{2}{K44}{1404}{}{如果一個社會對少數者權利根本沒有任何支持，違憲審查這個制度就無法為該少數的權利\mbox{做任何事。}}{A practice of judicial review cannot do anything for the rights of the minority if there is no support at all in the society for minority rights.}{社會對少數者權利全無支持時，違憲審查\mbox{無能為力。}}
\keyentry{2}{K144}{1405}{}{相反地，民選機關可能更善於保護少數者權利，因為選舉安排會提供一個管道，把民眾對少數者權利的支持導入立法機關，而通往司法機關並沒有這\mbox{樣的管道。}}{On the contrary, elective institutions may be better at protecting minority rights because electoral arrangements will provide a way of channeling popular support for minority rights into the legislature, whereas there are no such channels into the judiciary.}{對「同情少數者的是菁英」這一假定的回應：假定不成立時，民選機關保護少數者較有利；作者並指此假定合於關於民主的\mbox{古老偏見。}}
\keyentry{3}{K179}{1405}{}{他們會把這當作一項經驗性主張來辯護，但我必須說，它與關於民主決定的古老偏見\mbox{完全吻合。}}{They will defend this as an empirical claim, but I must say it is entirely consonant with ancient prejudices about democratic decisionmaking.}{「同情只見於菁英」之說，作者指為對民主決定的\mbox{古老偏見。}}
\keyentry{3}{K145}{1405}{}{但在制度上，司法的關照可能使事情變得更糟，或至少未能使事情\mbox{好轉多少。}}{But institutionally, judicial solicitude may make things worse, or at least fail to make them much better.}{長期而言違憲審查未必有助於矯正立法失能；1950、1960 年代的民權立法為例，關鍵在聯邦與州的分別。K8\mbox{1 承此。}}
\keyentry{3}{K180}{1405}{}{一如美國在 1950 與 1960 年代所發現的，儘管 Brown 案及其他案件中司法對種族隔離的抨擊令人振奮，最終所需要的是強力的立法介入（以 Civil Rights Act 的形式），而結果顯示，主要的差別並不在法院與立法機關之別本身，而在聯邦機關與州機關之別，最後由聯邦立法機關扮演決定\mbox{性的角色。}}{As the United States found in the 1950s and 1960s, for all the excitement of judicial attacks on segregation in Brown and other cases, what was needed in the end was strong legislative intervention (in the form of the Civil Rights Act), and it turned out that the main difference was not courts versus legislatures per se, but federal institutions versus state institutions, with the federal legislature finally playing the decisive role.}{長期效果：民權的決定性進展來自\mbox{國會立法。}}
\keyentry{2}{K181}{1405}{}{整體而言，我們不應把 Carolene Products 案的腳註或任何類似的法理，當作槓桿，藉以「撬出」一套更一般的違憲\mbox{審查制度。}}{Overall, we should not read the Carolene Products footnote or any similar doctrine as a way of ``leveraging'' a more general practice of judicial review into existence.}{第柒部分對腳註四的結論：它只支持特殊情形下的介入，不支持一般的違憲\mbox{審查制度。}}
\keyentry{4}{K198}{1406}{四級・生動表述}{孤立而隔絕的少數這個問題，不應被看作違憲審查的某種特洛伊木馬，也不應被拿來為難反對違憲審\mbox{查的論證。}}{The problem of discrete and insular minorities is not to be seen as a sort of Trojan Horse for judicial review or as a basis for embarrassing the arguments against it.}{特洛伊木馬之喻：腳註四不能用來夾帶一般的違憲審查；接 K181\mbox{、K45。}}
\keyentry{3}{K45}{1406}{}{考慮這類情形的目的不在於為違憲審查辯護；毋寧在於做那些最能保障受影響少數之\mbox{權利的事。}}{The aim of considering such cases is not to defend judicial review; rather it is to do whatever best secures the rights of the minorities affected.}{考慮非核心情形的目的在保障受影響少數\mbox{者的權利。}}
\keyentry{4}{K253}{1406}{四級・生動表述}{我們應當直接以此為目標，並且要意識到：並不存在令人信服而支持違憲審查的一般性論證，可以讓這個問題充當其意識\mbox{形態先鋒。}}{We should aim directly at that, conscious of the fact that there is no convincing general argument for judicial review of which this could be treated as a sort of ideological vanguard.}{「意識形態先鋒」的軍事比喻，與前句 K198 的特洛伊\mbox{木馬連讀。}}
\subsection*{結論}
\keyentry{1}{K46}{1406}{一級・結論一}{我所試圖顯示的是：一個相當民主的社會，若其主要問題不在立法制度失能，而在成員對權利有分歧，以權利為本的違憲審查為什麼不適合這\mbox{樣的社會。}}{Instead I have tried to show why rights-based judicial review is inappropriate for reasonably democratic societies whose main problem is not that their legislative institutions are dysfunctional but that their members disagree about rights.}{重述適用範圍，第貳部分的四項假設在結論裡收成「相當民主而成員對權利有分\mbox{歧」一語。}}
\keyentry{1}{K47}{1406}{一級・結論二}{我已論稱，普通的立法程序能夠做到這一點，而在其上再加一層由法院所為的終局審查，對這個程序幾乎毫無增益，只是以相當侮辱人的方式剝奪了人民的參政權，並以法條主義遮蔽了我們在權利分歧中所爭的\mbox{道德問題。}}{Ordinary legislative procedures can do this, I have argued, and an additional layer of final review by courts adds little to the process except a rather insulting form of disenfranchisement and a legalistic obfuscation of the moral issues at stake in our disagreements about rights.}{兩條論證在結論中各以一語重述：「剝奪參政權」承程序面，「把道德爭點弄得晦暗」\mbox{承結果面。}}
\keyentry{2}{K146}{1406}{}{但違憲審查的辯護者應當開始坦白地以此為據，主張這個制度，並在提出主張時，對促成這個制度的那些情況懷有幾分謙卑與羞愧，而不是四處宣講，把它奉為尊重權利的典範、現代立憲民主中正常且在規範上可\mbox{欲的要素。}}{But defenders of judicial review ought to start making their claims for the practice frankly on that basis—and make it with a degree of humility and shame in regard to the circumstances that elicit it—rather than preaching it abroad as the epitome of respect for rights and as a normal and normatively desirable element of modern constitutional democracy.}{全文末句：違憲審查只能以非核心情形的弊病為據主張，不能作為現代立憲民主的常態要素推廣。承 K42\mbox{、K47。}}